% !TEX root = ../main.tex
% 本文件由 build-tex.sh 从 35-DSP指南-索引中译.md 生成，改内容请改 md 后重新生成

\chapter{索引}\label{ch-35}

\section{A}

\begin{itemize}
\tightlist
\item
  \textbf{A-law companding} A 律压扩 --- 362-364 \dspsym{→}
  \hyperref[ch-22]{第 22 章}
\item
  \textbf{Accuracy} 准确度 --- 32-34 \dspsym{→} \hyperref[ch-02]{第 2
  章}
\item
  \textbf{Additivity} 可加性 --- 89-91 \dspsym{→} \hyperref[ch-05]{第 5
  章}；185-187 \dspsym{→} \hyperref[ch-10]{第 10 章}
\item
  \textbf{Algebraic reconstruction technique (ART)} 代数重建技术 ---
  444-445 \dspsym{→} \hyperref[ch-25]{第 25 章}
\item
  \textbf{Aliasing} 混叠

  \begin{itemize}
  \tightlist
  \item
    frequency domain：频域混叠 --- 196-200 \dspsym{→}
    \hyperref[ch-10]{第 10 章}；212-214, 220-222 \dspsym{→}
    \hyperref[ch-11]{第 11 章}；372 \dspsym{→} \hyperref[ch-22]{第 22
    章}
  \item
    in sampling：采样中的 --- 39-45 \dspsym{→} \hyperref[ch-03]{第 3 章}
  \item
    sinc function, equation for aliased：sinc 函数，混叠形式的公式 ---
    212-214 \dspsym{→} \hyperref[ch-11]{第 11 章}
  \item
    time domain：时域混叠 --- 194-196 \dspsym{→} \hyperref[ch-10]{第 10
    章}；300 \dspsym{→} \hyperref[ch-17]{第 17 章}
  \end{itemize}
\item
  \textbf{Alternating current (AC), defined} 交流（定义） --- 14
  \dspsym{→} \hyperref[ch-02]{第 2 章}
\item
  \textbf{Amplitude modulation (AM)} 幅度调制（调幅） --- 204-206
  \dspsym{→} \hyperref[ch-10]{第 10 章}；216-217 \dspsym{→}
  \hyperref[ch-11]{第 11 章}；370 \dspsym{→} \hyperref[ch-22]{第 22 章}
\item
  \textbf{Analysis equations} 分析方程 --- 见 \textbf{Fourier
  transform}（傅里叶变换）条下
\item
  \textbf{Antialias filter} 抗混叠滤波器 --- 见 \textbf{filters-
  analog}（滤波器，模拟）条下
\item
  \textbf{Arithmetic encoding} 算术编码 --- 486 \dspsym{→}
  \hyperref[ch-27]{第 27 章}
\item
  \textbf{Artificial neural net} 人工神经网络 --- 458 \dspsym{→}
  \hyperref[ch-26]{第 26 章}，另见 \textbf{Neural network}（神经网络）
\item
  \textbf{Artificial reverberation in music} 音乐中的人工混响 --- 5
  \dspsym{→} \hyperref[ch-01]{第 1 章}
\item
  \textbf{ASCII codes, table of} ASCII 码表 --- 484 \dspsym{→}
  \hyperref[ch-27]{第 27 章}
\item
  \textbf{Aspect ratio of television} 电视的宽高比 --- 386 \dspsym{→}
  \hyperref[ch-23]{第 23 章}
\item
  \textbf{Assembly program} 汇编程序 --- 76-77 \dspsym{→}
  \hyperref[ch-04]{第 4 章}；520 \dspsym{→} \hyperref[ch-28]{第 28 章}
\item
  \textbf{Astrophotography} 天文摄影 --- 1, 10 \dspsym{→}
  \hyperref[ch-01]{第 1 章}；373-375, 394-396 \dspsym{→}
  \hyperref[ch-23]{第 23 章}
\item
  \textbf{Audio processing} 音频处理 --- 5-7 \dspsym{→}
  \hyperref[ch-01]{第 1 章}；304-307 \dspsym{→} \hyperref[ch-17]{第 17
  章}；311 \dspsym{→} \hyperref[ch-18]{第 18 章}；351-372 \dspsym{→}
  \hyperref[ch-22]{第 22 章}
\item
  \textbf{Audio signaling tones, detection of} 音频信令音的检测 --- 293
  \dspsym{→} \hyperref[ch-16]{第 16 章}
\item
  \textbf{Automatic gain control (AGC)} 自动增益控制 --- 370 \dspsym{→}
  \hyperref[ch-22]{第 22 章}
\end{itemize}

\section{B}

\begin{itemize}
\tightlist
\item
  \textbf{Backprojection} 反投影 --- 446-450 \dspsym{→}
  \hyperref[ch-25]{第 25 章}
\item
  \textbf{Basis functions} 基函数

  \begin{itemize}
  \tightlist
  \item
    discrete Fourier transform：离散傅里叶变换的基函数 --- 150-152,
    158-159 \dspsym{→} \hyperref[ch-08]{第 8 章}
  \item
    discrete cosine transform：离散余弦变换的基函数 --- 496-497
    \dspsym{→} \hyperref[ch-27]{第 27 章}
  \end{itemize}
\item
  \textbf{Bessel filter} 贝塞尔滤波器 --- 见 \textbf{Filters,
  analog}（滤波器，模拟）
\item
  \textbf{Bias node in neural networks} 神经网络中的偏置节点 --- 462-463
  \dspsym{→} \hyperref[ch-26]{第 26 章}
\item
  \textbf{Bilinear interpolation} 双线性插值 --- 396 \dspsym{→}
  \hyperref[ch-23]{第 23 章}
\item
  \textbf{Binary image processing} 二值图像处理 --- 见
  \textbf{Morphological processing}（形态学处理）
\item
  \textbf{Biquad} 双二阶节 --- 600 \dspsym{→} \hyperref[ch-32]{第 32 章}
\item
  \textbf{Bit map to vector map conversion} 位图到矢量图的转换 --- 442
  \dspsym{→} \hyperref[ch-25]{第 25 章}
\item
  \textbf{Bit reversal sorting in FFT} FFT 中的位反转排序 --- 229
  \dspsym{→} \hyperref[ch-12]{第 12 章}
\item
  \textbf{Blacker than black video} 视频中的「比黑更黑」电平 --- 385
  \dspsym{→} \hyperref[ch-23]{第 23 章}
\item
  \textbf{Blob analysis} 斑点分析 --- 436 \dspsym{→} \hyperref[ch-25]{第
  25 章}，另见 \textbf{Morphological processing}（形态学处理）
\item
  \textbf{Brackets, indicating discrete signals} 方括号，表示离散信号
  --- 87 \dspsym{→} \hyperref[ch-05]{第 5 章}
\item
  \textbf{Brightness in images} 图像的亮度 --- 387-391 \dspsym{→}
  \hyperref[ch-23]{第 23 章}
\item
  \textbf{Butterfly calculation in FFT} FFT 中的蝶形运算 --- 231-232
  \dspsym{→} \hyperref[ch-12]{第 12 章}
\item
  \textbf{Butterworth filter} 巴特沃斯滤波器 --- 见
  \textbf{Filters}（滤波器）条下
\end{itemize}

\section{C}

\begin{itemize}
\tightlist
\item
  \textbf{C program} C 程序 --- 67, 77 \dspsym{→} \hyperref[ch-04]{第 4
  章}；520 \dspsym{→} \hyperref[ch-28]{第 28 章}
\item
  \textbf{Cascaded stages} 级联 --- 96 \dspsym{→} \hyperref[ch-05]{第 5
  章}；133 \dspsym{→} \hyperref[ch-07]{第 7 章}，另见 \textbf{Filters-
  recursive}（滤波器，递归）条下
\item
  \textbf{Caruso, restoration of recordings by} 卡鲁索录音的修复 ---
  304-307 \dspsym{→} \hyperref[ch-17]{第 17 章}
\item
  \textbf{CAT scanner} CAT 扫描仪 --- 见 \textbf{Computed
  tomography}（计算机断层成像）
\item
  \textbf{Causal signals and systems} 因果信号与因果系统 --- 130
  \dspsym{→} \hyperref[ch-07]{第 7 章}
\item
  \textbf{CCD} CCD --- 见 \textbf{Charge coupled device}（电荷耦合器件）
\item
  \textbf{Central limit theorem} 中心极限定理 --- 30 \dspsym{→}
  \hyperref[ch-02]{第 2 章}；135-136 \dspsym{→} \hyperref[ch-07]{第 7
  章}；407 \dspsym{→} \hyperref[ch-24]{第 24 章}
\item
  \textbf{Cepstrum} 倒谱 --- 371 \dspsym{→} \hyperref[ch-22]{第 22 章}
\item
  \textbf{Charge coupled device (CCD)} 电荷耦合器件 --- 381-385
  \dspsym{→} \hyperref[ch-23]{第 23 章}；430-432 \dspsym{→}
  \hyperref[ch-25]{第 25 章}
\item
  \textbf{Charge sensitive amplifier, in CCD} 电荷敏感放大器（CCD 中）
  --- 382-384 \dspsym{→} \hyperref[ch-23]{第 23 章}
\item
  \textbf{Chebyshev filter} 切比雪夫滤波器 --- 见
  \textbf{Filters}（滤波器）条下
\item
  \textbf{Chirp signals and systems} 啁啾信号与啁啾系统 --- 222-224
  \dspsym{→} \hyperref[ch-11]{第 11 章}
\item
  \textbf{Chrominance signal, in television} 色度信号（电视中） --- 386
  \dspsym{→} \hyperref[ch-23]{第 23 章}
\item
  \textbf{Circular buffer} 循环缓冲 --- 507 \dspsym{→}
  \hyperref[ch-28]{第 28 章}
\item
  \textbf{Circularity} 循环性 --- 见 \textbf{Discrete Fourier
  transform}（离散傅里叶变换）条下
\item
  \textbf{Classifiers} 分类器 --- 458 \dspsym{→} \hyperref[ch-26]{第 26
  章}
\item
  \textbf{Close neighbors in images} 图像中的近邻 --- 439 \dspsym{→}
  \hyperref[ch-25]{第 25 章}
\item
  \textbf{Closing, morphological} 闭运算（形态学） --- 437 \dspsym{→}
  \hyperref[ch-25]{第 25 章}
\item
  \textbf{Coefficient of variation (CV)} 变异系数 --- 17 \dspsym{→}
  \hyperref[ch-02]{第 2 章}
\item
  \textbf{Color} 颜色 --- 376, 379-381, 386 \dspsym{→}
  \hyperref[ch-23]{第 23 章}
\item
  \textbf{Compact laser disc (CD)} 激光唱片（CD） --- 359-362 \dspsym{→}
  \hyperref[ch-22]{第 22 章}
\item
  \textbf{Complex logarithm} 复对数 --- 372 \dspsym{→}
  \hyperref[ch-22]{第 22 章}
\item
  \textbf{Complex numbers} 复数

  \begin{itemize}
  \tightlist
  \item
    addition：加法 --- 553-554 \dspsym{→} \hyperref[ch-30]{第 30 章}
  \item
    associative property：结合律 --- 554 \dspsym{→} \hyperref[ch-30]{第
    30 章}
  \item
    commutative property：交换律 --- 5054 \dspsym{→} 索引
  \item
    complex number system：复数系统 --- 551-554 \dspsym{→}
    \hyperref[ch-30]{第 30 章}
  \item
    complex plane：复平面 --- 552-553 \dspsym{→} \hyperref[ch-30]{第 30
    章}
  \item
    conjugation：共轭 --- 193-194 \dspsym{→} \hyperref[ch-10]{第 10 章}
  \item
    distributive property：分配律 --- 554 \dspsym{→} \hyperref[ch-30]{第
    30 章}
  \item
    division：除法 --- 554, 557 \dspsym{→} \hyperref[ch-30]{第 30 章}
  \item
    Euler's relation：欧拉关系 --- 556 \dspsym{→} \hyperref[ch-30]{第 30
    章}；569-570 \dspsym{→} \hyperref[ch-31]{第 31 章}
  \item
    exponential form：指数形式 --- 557 \dspsym{→} \hyperref[ch-30]{第 30
    章}；569-570 \dspsym{→} \hyperref[ch-31]{第 31 章}；584 \dspsym{→}
    \hyperref[ch-32]{第 32 章}
  \item
    multiplication：乘法 --- 554, 557 \dspsym{→} \hyperref[ch-30]{第 30
    章}
  \item
    polar notation：极坐标表示法 --- 555-557 \dspsym{→}
    \hyperref[ch-30]{第 30 章}
  \item
    rectangular notation：直角坐标表示法 --- 552-556 \dspsym{→}
    \hyperref[ch-30]{第 30 章}
  \item
    sinusoids, representing：表示正弦波 --- 559-561 \dspsym{→}
    \hyperref[ch-30]{第 30 章}
  \item
    subtraction：减法 --- 554 \dspsym{→} \hyperref[ch-30]{第 30 章}
  \item
    systems, representing：表示系统 --- 561-563 \dspsym{→}
    \hyperref[ch-30]{第 30 章}
  \end{itemize}
\item
  \textbf{Companding} 压扩 --- 4 \dspsym{→} \hyperref[ch-01]{第 1
  章}；358-359, 362-364 \dspsym{→} \hyperref[ch-22]{第 22 章}
\item
  \textbf{Compiler} 编译器 --- 78 \dspsym{→} \hyperref[ch-04]{第 4
  章}；546 \dspsym{→} \hyperref[ch-29]{第 29 章}
\item
  \textbf{Composite video} 复合视频 --- 384-386 \dspsym{→}
  \hyperref[ch-23]{第 23 章}
\item
  \textbf{Computed tomography (CT)} 计算机断层成像 --- 9 \dspsym{→}
  \hyperref[ch-01]{第 1 章}；411 \dspsym{→} \hyperref[ch-24]{第 24
  章}；429, 442-450 \dspsym{→} \hyperref[ch-25]{第 25 章}
\item
  \textbf{Compression, data} 压缩，数据 --- 见 \textbf{Data
  compression}（数据压缩）
\item
  \textbf{Compression \& expansion of signals} 信号的压缩与扩展 ---
  200-203 \dspsym{→} \hyperref[ch-10]{第 10 章}
\item
  \textbf{Compression ratio in JPEG} JPEG 的压缩比 --- 500 \dspsym{→}
  \hyperref[ch-27]{第 27 章}
\item
  \textbf{Connectivity analysis} 连通性分析 --- 434 \dspsym{→}
  \hyperref[ch-25]{第 25 章}，另见 \textbf{Morphological
  processing}（形态学处理）
\item
  \textbf{Continuous signal (defined)} 连续信号（定义） --- 11
  \dspsym{→} \hyperref[ch-02]{第 2 章}
\item
  \textbf{Contrast, image} 对比度，图像 --- 387-391 \dspsym{→}
  \hyperref[ch-23]{第 23 章}；432-434 \dspsym{→} \hyperref[ch-25]{第 25
  章}
\item
  \textbf{Convolution} 卷积

  \begin{itemize}
  \tightlist
  \item
    associative property：结合律 --- 133 \dspsym{→} \hyperref[ch-07]{第
    7 章}
  \item
    circular：循环卷积 --- 182-184 \dspsym{→} \hyperref[ch-09]{第 9
    章}；314 \dspsym{→} \hyperref[ch-18]{第 18 章}
  \item
    commutative property：交换律 --- 113 \dspsym{→} \hyperref[ch-06]{第
    6 章}；132 \dspsym{→} \hyperref[ch-07]{第 7 章}
  \item
    continuous：连续卷积 --- 246-252 \dspsym{→} \hyperref[ch-13]{第 13
    章}
  \item
    convolution machine：卷积机 --- 116-119 \dspsym{→}
    \hyperref[ch-06]{第 6 章}
  \item
    discrete：离散卷积 --- 107-122 \dspsym{→} \hyperref[ch-06]{第 6 章}
  \item
    distributive property：分配律 --- 134 \dspsym{→} \hyperref[ch-07]{第
    7 章}
  \item
    end effects：端点效应 --- 118-121 \dspsym{→} \hyperref[ch-06]{第 6
    章}；408 \dspsym{→} \hyperref[ch-24]{第 24 章}
  \item
    execution time：执行时间 --- 140 \dspsym{→} \hyperref[ch-07]{第 7
    章}；316 \dspsym{→} \hyperref[ch-18]{第 18 章}；318-422 \dspsym{→}
    \hyperref[ch-18]{第 18 章}、\hyperref[ch-24]{第 24 章}
  \item
    frequency multiplication, by：通过频域相乘实现 --- 179-184
    \dspsym{→} \hyperref[ch-09]{第 9 章}，另见 \textbf{FFT
    convolution}（FFT 卷积）
  \item
    image：图像卷积 --- 397-410, 416-422 \dspsym{→} \hyperref[ch-24]{第
    24 章}
  \item
    immersion of impulse response：冲激响应的浸没 --- 119 \dspsym{→}
    \hyperref[ch-06]{第 6 章}；410 \dspsym{→} \hyperref[ch-24]{第 24 章}
  \item
    input side view：输入端视角 --- 112-115 \dspsym{→}
    \hyperref[ch-06]{第 6 章}；246-247 \dspsym{→} \hyperref[ch-13]{第 13
    章}；398 \dspsym{→} \hyperref[ch-24]{第 24 章}；418-423 \dspsym{→}
    \hyperref[ch-24]{第 24 章}、\hyperref[ch-25]{第 25 章}
  \item
    left-for-right flip：左右翻转 --- 117 \dspsym{→} \hyperref[ch-06]{第
    6 章}；138-140 \dspsym{→} \hyperref[ch-07]{第 7 章}；194 \dspsym{→}
    \hyperref[ch-10]{第 10 章}；416-421 \dspsym{→} \hyperref[ch-24]{第
    24 章}
  \item
    neural networks, carried out by：用神经网络实现 --- 464-465
    \dspsym{→} \hyperref[ch-26]{第 26 章}
  \item
    output side view：输出端视角 --- 116-121 \dspsym{→}
    \hyperref[ch-06]{第 6 章}；246-247 \dspsym{→} \hyperref[ch-13]{第 13
    章}；399-400, 418-421 \dspsym{→} \hyperref[ch-24]{第 24 章}
  \item
    piecewise polynomial method：分段多项式法 --- 250-252 \dspsym{→}
    \hyperref[ch-13]{第 13 章}
  \item
    separability, image convolution by：可分离性，用于图像卷积 ---
    404-407 \dspsym{→} \hyperref[ch-24]{第 24 章}
  \item
    sum of weighted inputs, viewed as：看作输入的加权和 --- 122
    \dspsym{→} \hyperref[ch-06]{第 6 章}
  \end{itemize}
\item
  \textbf{Cooley, J.W. \& Tukey, J.W.} 库利与图基 --- 225 \dspsym{→}
  \hyperref[ch-12]{第 12 章}
\item
  \textbf{Correlation} 相关 --- 136-140 \dspsym{→} \hyperref[ch-07]{第 7
  章}

  \begin{itemize}
  \tightlist
  \item
    autocorrelation：自相关 --- 137 \dspsym{→} \hyperref[ch-07]{第 7 章}
  \item
    convolution, carried out by：用卷积实现 --- 138-140 \dspsym{→}
    \hyperref[ch-07]{第 7 章}；194 \dspsym{→} \hyperref[ch-10]{第 10
    章}；416-422 \dspsym{→} \hyperref[ch-24]{第 24 章}
  \item
    correlation machine：相关机 --- 138-140 \dspsym{→}
    \hyperref[ch-07]{第 7 章}
  \item
    cross-correlation：互相关 --- 137 \dspsym{→} \hyperref[ch-07]{第 7
    章}
  \item
    Fourier transform use of：在傅里叶变换中的应用 --- 见
    \textbf{Discrete Fourier transform}（离散傅里叶变换）
  \item
    matched filter：匹配滤波器 --- 见 \textbf{Filters}（滤波器）
  \item
    neural networks, carried out by：用神经网络实现 --- 464 \dspsym{→}
    \hyperref[ch-26]{第 26 章}
  \item
    radar and sonar use：在雷达与声呐中的应用 --- 137 \dspsym{→}
    \hyperref[ch-07]{第 7 章}
  \end{itemize}
\item
  \textbf{Counting statistics} 计数统计 --- 434-436 \dspsym{→}
  \hyperref[ch-25]{第 25 章}
\item
  \textbf{CRT display, point spread function of} CRT 显示器的点扩散函数
  --- 424 \dspsym{→} \hyperref[ch-25]{第 25 章}
\item
  \textbf{Cumulative distribution function} 累积分布函数 --- 28-29
  \dspsym{→} \hyperref[ch-02]{第 2 章}
\item
  \textbf{CVSD modulation} CVSD 调制（连续可变斜率增量调制） --- 62-64
  \dspsym{→} \hyperref[ch-03]{第 3 章}
\end{itemize}

\section{D}

\begin{itemize}
\tightlist
\item
  \textbf{Data compression} 数据压缩 --- 4, 10 \dspsym{→}
  \hyperref[ch-01]{第 1 章}；481-502 \dspsym{→} \hyperref[ch-27]{第 27
  章}
\item
  \textbf{dB} dB --- 见 \textbf{Decibels}（分贝）
\item
  \textbf{dBm} dBm --- 264 \dspsym{→} \hyperref[ch-14]{第 14 章}
\item
  \textbf{dBV} dBV --- 264 \dspsym{→} \hyperref[ch-14]{第 14 章}
\item
  \textbf{DC offset} 直流偏置

  \begin{itemize}
  \tightlist
  \item
    in Fourier transform：傅里叶变换中的 --- 152 \dspsym{→}
    \hyperref[ch-08]{第 8 章}
  \item
    linearity of：直流偏置的线性 --- 97 \dspsym{→} \hyperref[ch-05]{第 5
    章}
  \end{itemize}
\item
  \textbf{DCT} DCT --- 见 \textbf{Discrete cosine
  transform}（离散余弦变换）
\item
  \textbf{Decibels} 分贝 --- 263-264 \dspsym{→} \hyperref[ch-14]{第 14
  章}
\item
  \textbf{Decibels SPL} 声压级分贝 --- 352-356 \dspsym{→}
  \hyperref[ch-22]{第 22 章}
\item
  \textbf{Decimation} 抽取 --- 见 \textbf{Multirate}
\item
  \textbf{Decimation in frequency FFT} 按频率抽取 FFT --- 234 \dspsym{→}
  \hyperref[ch-12]{第 12 章}
\item
  \textbf{Decimation in time FFT} 按时间抽取 FFT --- 234 \dspsym{→}
  \hyperref[ch-12]{第 12 章}
\item
  \textbf{Decomposition} 分解

  \begin{itemize}
  \tightlist
  \item
    defined：定义 --- 98 \dspsym{→} \hyperref[ch-05]{第 5 章}
  \item
    even/odd decomposition：偶/奇分解 --- 102-103 \dspsym{→}
    \hyperref[ch-05]{第 5 章}；240-242 \dspsym{→} \hyperref[ch-12]{第 12
    章}
  \item
    Fourier decomposition：傅里叶分解 --- 104, 105 \dspsym{→}
    \hyperref[ch-05]{第 5 章}；147 \dspsym{→} \hyperref[ch-08]{第 8 章}
  \item
    impulse decomposition：冲激分解 --- 100-101 \dspsym{→}
    \hyperref[ch-05]{第 5 章}
  \item
    interlaced decomposition：交错分解 --- 103-104 \dspsym{→}
    \hyperref[ch-05]{第 5 章}；228-229 \dspsym{→} \hyperref[ch-12]{第 12
    章}
  \item
    step decomposition：阶跃分解 --- 100-101 \dspsym{→}
    \hyperref[ch-05]{第 5 章}
  \item
    strategy for using：使用策略 --- 98-100 \dspsym{→}
    \hyperref[ch-05]{第 5 章}
  \end{itemize}
\item
  \textbf{Deconvolution} 反卷积 --- 179-180 \dspsym{→}
  \hyperref[ch-09]{第 9 章}；300-307 \dspsym{→} \hyperref[ch-17]{第 17
  章}
\item
  \textbf{Delta encoding, data compression} 差分编码，数据压缩 ---
  486-488 \dspsym{→} \hyperref[ch-27]{第 27 章}
\item
  \textbf{Delta function} 冲激函数

  \begin{itemize}
  \tightlist
  \item
    continuous：连续冲激函数 --- 243-245 \dspsym{→} \hyperref[ch-13]{第
    13 章}
  \item
    discrete：离散冲激函数 --- 107-109 \dspsym{→} \hyperref[ch-06]{第 6
    章}
  \item
    Fourier transform of：其傅里叶变换 --- 200 \dspsym{→}
    \hyperref[ch-10]{第 10 章}；209-212 \dspsym{→} \hyperref[ch-11]{第
    11 章}
  \item
    identity for convolution：卷积的单位元 --- 123 \dspsym{→}
    \hyperref[ch-07]{第 7 章}
  \item
    two-dimensional (image)：二维（图像） --- 397-398 \dspsym{→}
    \hyperref[ch-24]{第 24 章}
  \end{itemize}
\item
  \textbf{Delta modulation} 增量调制 --- 60-66 \dspsym{→}
  \hyperref[ch-03]{第 3 章}
\item
  \textbf{Delta-sigma} 增量-累加（Δ-Σ） --- 63-66 \dspsym{→}
  \hyperref[ch-03]{第 3 章}
\item
  \textbf{Dependent variable (defined)} 因变量（定义） --- 12 \dspsym{→}
  \hyperref[ch-02]{第 2 章}
\item
  \textbf{Derivative, discrete} 导数，离散 --- 见 \textbf{First
  difference}（一阶差分）
\item
  \textbf{Derivative of continuous step function} 连续阶跃函数的导数 ---
  250-252 \dspsym{→} \hyperref[ch-13]{第 13 章}
\item
  \textbf{DFT} DFT --- 见 \textbf{Discrete Fourier
  transform}（离散傅里叶变换）
\item
  \textbf{Differential equations, RLC circuits} 微分方程，RLC 电路 ---
  563-564 \dspsym{→} \hyperref[ch-30]{第 30 章}
\item
  \textbf{Difference equation} 差分方程 --- 见 \textbf{Recursion
  equation}（递推方程）
\item
  \textbf{Digital-to-analog conversion (DAC)} 数模转换 --- 44-48
  \dspsym{→} \hyperref[ch-03]{第 3 章}
\item
  \textbf{Dilatation} 膨胀（异体拼写） --- 见 \textbf{Dilation}（膨胀）
\item
  \textbf{Dilation} 膨胀 --- 437 \dspsym{→} \hyperref[ch-25]{第 25
  章}，另见 \textbf{Morphological processing}（形态学处理）
\item
  \textbf{Direct current (DC) (defined)} 直流（定义） --- 14 \dspsym{→}
  \hyperref[ch-02]{第 2 章}
\item
  \textbf{Discrete signal (defined)} 离散信号（定义） --- 11 \dspsym{→}
  \hyperref[ch-02]{第 2 章}
\item
  \textbf{Discrete cosine transform (DCT)} 离散余弦变换 --- 496
  \dspsym{→} \hyperref[ch-27]{第 27 章}
\item
  \textbf{Discrete Fourier transform (DFT)} 离散傅里叶变换

  \begin{itemize}
  \tightlist
  \item
    另见 \textbf{Fourier transform}（傅里叶变换）；\textbf{Fourier
    transform pairs}（傅里叶变换对）
  \item
    analysis：分析 --- 157-161 \dspsym{→} \hyperref[ch-08]{第 8
    章}；567-580 \dspsym{→} \hyperref[ch-31]{第 31 章}
  \item
    basis functions：基函数 --- 150-152, 158-159 \dspsym{→}
    \hyperref[ch-08]{第 8 章}
  \item
    circular：循环性 --- 195 \dspsym{→} \hyperref[ch-10]{第 10 章}，另见
    \textbf{periodic}
  \item
    complex DFT：复数 DFT --- 225-227 \dspsym{→} \hyperref[ch-12]{第 12
    章}；519-532 \dspsym{→} \hyperref[ch-28]{第 28
    章}（疑为旧版页码，按内容应在 567-580 \dspsym{→} \hyperref[ch-31]{第
    31 章}）
  \item
    correlation method：相关法 --- 157-160 \dspsym{→}
    \hyperref[ch-08]{第 8 章}
  \item
    examples of：示例 --- 142-143 \dspsym{→} \hyperref[ch-08]{第 8
    章}；171, 181 \dspsym{→} \hyperref[ch-09]{第 9 章}；186, 193
    \dspsym{→} \hyperref[ch-10]{第 10 章}，另见 \textbf{Fourier
    Transform pairs}（傅里叶变换对）
  \item
    forward DFT：正向 DFT --- 157-160 \dspsym{→} \hyperref[ch-08]{第 8
    章}；567-580 \dspsym{→} \hyperref[ch-31]{第 31 章}
  \item
    frequency resolution：频率分辨率 --- 173-176 \dspsym{→}
    \hyperref[ch-09]{第 9 章}
  \item
    inverse DFT：逆 DFT --- 152-156 \dspsym{→} \hyperref[ch-08]{第 8
    章}；573-580 \dspsym{→} \hyperref[ch-31]{第 31 章}
  \item
    negative frequencies：负频率 --- 196-200 \dspsym{→}
    \hyperref[ch-10]{第 10 章}；209-211 \dspsym{→} \hyperref[ch-11]{第
    11 章}；226-227 \dspsym{→} \hyperref[ch-12]{第 12 章}；569-575
    \dspsym{→} \hyperref[ch-31]{第 31 章}
  \item
    orthognal basis functions：正交基函数 --- 158-159 \dspsym{→}
    \hyperref[ch-08]{第 8 章}
  \item
    periodic frequency domain：周期性的频域 --- 196-200 \dspsym{→}
    \hyperref[ch-10]{第 10 章}
  \item
    periodic time domain：周期性的时域 --- 194-196 \dspsym{→}
    \hyperref[ch-10]{第 10 章}
  \item
    real DFT：实数 DFT --- 141-160 \dspsym{→} \hyperref[ch-08]{第 8
    章}；225-227 \dspsym{→} \hyperref[ch-12]{第 12 章}；519-521
    \dspsym{→} \hyperref[ch-28]{第 28 章}（疑为旧版页码，按内容应在
    567-569 \dspsym{→} \hyperref[ch-31]{第 31 章}）
  \item
    spectral density：谱密度 --- 156 \dspsym{→} \hyperref[ch-08]{第 8
    章}
  \item
    spectral leakage：频谱泄漏 --- 175-176 \dspsym{→}
    \hyperref[ch-09]{第 9 章}
  \item
    synthesis：合成 --- 152-156 \dspsym{→} \hyperref[ch-08]{第 8
    章}；573-580 \dspsym{→} \hyperref[ch-31]{第 31 章}
  \end{itemize}
\item
  \textbf{Discrete Time Fourier Transform (DTFT)} 离散时间傅里叶变换 ---
  见 \textbf{Fourier transform}（傅里叶变换）
\item
  \textbf{Distant neighbors in images} 图像中的远邻 --- 439 \dspsym{→}
  \hyperref[ch-25]{第 25 章}
\item
  \textbf{Dithering} 抖动 --- 38-39 \dspsym{→} \hyperref[ch-03]{第 3
  章}；374 \dspsym{→} \hyperref[ch-23]{第 23 章}
\item
  \textbf{Division of frequency domain signals} 频域信号的除法 ---
  181-182 \dspsym{→} \hyperref[ch-09]{第 9 章}
\item
  \textbf{DN (digital number) in images} 图像中的 DN（数字值） --- 374
  \dspsym{→} \hyperref[ch-23]{第 23 章}
\item
  \textbf{Dolby stereo} 杜比立体声 --- 362 \dspsym{→}
  \hyperref[ch-22]{第 22 章}
\item
  \textbf{Domain (defined)} 域（定义） --- 12 \dspsym{→}
  \hyperref[ch-02]{第 2 章}
\item
  \textbf{Double precision} 双精度 --- 70-74 \dspsym{→}
  \hyperref[ch-04]{第 4 章}；284 \dspsym{→} \hyperref[ch-15]{第 15
  章}；339 \dspsym{→} \hyperref[ch-20]{第 20 章}
\item
  \textbf{DSP microprocessor} DSP 微处理器 --- 84 \dspsym{→}
  \hyperref[ch-04]{第 4 章}
\item
  \textbf{DTFT} DTFT --- 见 \textbf{Fourier transform}（傅里叶变换）
\item
  \textbf{Duality} 对偶性 --- 161 \dspsym{→} \hyperref[ch-08]{第 8
  章}；210-212 \dspsym{→} \hyperref[ch-11]{第 11 章}；236 \dspsym{→}
  \hyperref[ch-12]{第 12 章}
\item
  \textbf{Dynamic range} 动态范围 --- 261-262 \dspsym{→}
  \hyperref[ch-14]{第 14 章}；378 \dspsym{→} \hyperref[ch-23]{第 23 章}
\end{itemize}

\section{E}

\begin{itemize}
\tightlist
\item
  \textbf{Ear} 耳朵 --- 351-355 \dspsym{→} \hyperref[ch-22]{第 22 章}
\item
  \textbf{Echo control in telephones} 电话中的回声控制 --- 5 \dspsym{→}
  \hyperref[ch-01]{第 1 章}
\item
  \textbf{Echo location} 回声定位 --- 7-8 \dspsym{→} \hyperref[ch-01]{第
  1 章}，另见
  \textbf{Radar}（雷达）；\textbf{Sonar}（声呐）；\textbf{Seismology}（地震学）
\item
  \textbf{Echoes in music} 音乐中的回声 --- 5 \dspsym{→}
  \hyperref[ch-01]{第 1 章}
\item
  \textbf{Edge detection} 边缘检测 --- 402-403, 416-417 \dspsym{→}
  \hyperref[ch-24]{第 24 章}
\item
  \textbf{Edge response} 边缘响应 --- 426-430 \dspsym{→}
  \hyperref[ch-25]{第 25 章}
\item
  \textbf{EFM (8 to 14 modulation)} EFM（8-14 调制） --- 360 \dspsym{→}
  \hyperref[ch-22]{第 22 章}
\item
  \textbf{Electric circuit analysis} 电路分析

  \begin{itemize}
  \tightlist
  \item
    phasor transform method：相量变换法 --- 563-566 \dspsym{→}
    \hyperref[ch-30]{第 30 章}
  \item
    Laplace transform method：拉普拉斯变换法 --- 592-599 \dspsym{→}
    \hyperref[ch-32]{第 32 章}
  \end{itemize}
\item
  \textbf{Electroencephalogram (EEG)} 脑电图 --- 292-293 \dspsym{→}
  \hyperref[ch-16]{第 16 章}
\item
  \textbf{End effect} 端点效应 --- 见 \textbf{Convolution}（卷积）条下
\item
  \textbf{Erosion} 腐蚀 --- 435 \dspsym{→} \hyperref[ch-25]{第 25
  章}，另见 \textbf{Morphological processing}（形态学处理）
\item
  \textbf{Euler's relation} 欧拉关系 --- 见 \textbf{Complex
  numbers}（复数）条下
\item
  \textbf{Even field in video} 视频中的偶场 --- 384-385 \dspsym{→}
  \hyperref[ch-23]{第 23 章}
\item
  \textbf{Even/odd decomposition} 偶/奇分解 --- 102 \dspsym{→}
  \hyperref[ch-05]{第 5 章}；240-242 \dspsym{→} \hyperref[ch-12]{第 12
  章}
\item
  \textbf{Even symmetry} 偶对称 --- 102 \dspsym{→} \hyperref[ch-05]{第 5
  章}；196 \dspsym{→} \hyperref[ch-10]{第 10 章}；240-241 \dspsym{→}
  \hyperref[ch-12]{第 12 章}
\item
  \textbf{Evolution, neural network learning} 进化，神经网络学习 --- 470
  \dspsym{→} \hyperref[ch-26]{第 26 章}
\item
  \textbf{Execution speed} 执行速度 --- 见 \textbf{Speed}（速度）
\item
  \textbf{Exponential, two ways to generate} 指数，两种生成方法 ---
  606-607 \dspsym{→} \hyperref[ch-33]{第 33 章}
\item
  \textbf{Eye} 眼睛 --- 376-381 \dspsym{→} \hyperref[ch-23]{第 23
  章}；404, 409 \dspsym{→} \hyperref[ch-24]{第 24 章}；432-434
  \dspsym{→} \hyperref[ch-25]{第 25 章}
\end{itemize}

\section{F}

\begin{itemize}
\tightlist
\item
  \textbf{False-positive (false-negative)} 假阳性（假阴性） --- 453-454
  \dspsym{→} \hyperref[ch-26]{第 26 章}
\item
  \textbf{Fast Fourier transform (FFT)} 快速傅里叶变换 --- 180
  \dspsym{→} \hyperref[ch-09]{第 9 章}；225-242 \dspsym{→}
  \hyperref[ch-12]{第 12 章}
\item
  \textbf{Feature extraction} 特征提取 --- 458 \dspsym{→}
  \hyperref[ch-26]{第 26 章}
\item
  \textbf{FFT convolution} FFT 卷积 --- 140 \dspsym{→}
  \hyperref[ch-07]{第 7 章}；179-184 \dspsym{→} \hyperref[ch-09]{第 9
  章}；311-318 \dspsym{→} \hyperref[ch-18]{第 18 章}；411, 416-422
  \dspsym{→} \hyperref[ch-24]{第 24 章}
\item
  \textbf{Field, television} 场，电视 --- 384-386 \dspsym{→}
  \hyperref[ch-23]{第 23 章}
\item
  \textbf{Filters, analog} 滤波器，模拟

  \begin{itemize}
  \tightlist
  \item
    另见 \textbf{Filters, digital}（滤波器，数字）；\textbf{Filters,
    recursive}（滤波器，递归）
  \item
    in ADC and DAC：在 ADC 与 DAC 中 --- 48-59 \dspsym{→}
    \hyperref[ch-03]{第 3 章}
  \item
    antialias：抗混叠 --- 48-49, 55-59 \dspsym{→} \hyperref[ch-03]{第 3
    章}；172-173 \dspsym{→} \hyperref[ch-09]{第 9 章}
  \item
    Bessel：贝塞尔 --- 49-59 \dspsym{→} \hyperref[ch-03]{第 3 章}；330
    \dspsym{→} \hyperref[ch-19]{第 19 章}；361 \dspsym{→}
    \hyperref[ch-22]{第 22 章}
  \item
    Butterworth：巴特沃斯 --- 49-58 \dspsym{→} \hyperref[ch-03]{第 3
    章}；600-604 \dspsym{→} \hyperref[ch-32]{第 32 章}
  \item
    Chebyshev：切比雪夫 --- 49-58 \dspsym{→} \hyperref[ch-03]{第 3
    章}；600-604 \dspsym{→} \hyperref[ch-32]{第 32 章}
  \item
    design methods：设计方法 --- 49-59 \dspsym{→} \hyperref[ch-03]{第 3
    章}；600-604 \dspsym{→} \hyperref[ch-32]{第 32 章}
  \item
    digital, compared to：与数字滤波器比较 --- 261-262 \dspsym{→}
    \hyperref[ch-14]{第 14 章}；343-345 \dspsym{→} \hyperref[ch-21]{第
    21 章}
  \item
    elliptic：椭圆 --- 54 \dspsym{→} \hyperref[ch-03]{第 3 章}
  \item
    frequency response：频率响应 --- 52-54 \dspsym{→}
    \hyperref[ch-03]{第 3 章}；176-179 \dspsym{→} \hyperref[ch-09]{第 9
    章}；598-599 \dspsym{→} \hyperref[ch-32]{第 32 章}
  \item
    high-pass：高通 --- 51 \dspsym{→} \hyperref[ch-03]{第 3 章}
  \item
    low-pass：低通 --- 49-54 \dspsym{→} \hyperref[ch-03]{第 3 章}；322
    \dspsym{→} \hyperref[ch-19]{第 19 章}；343-345 \dspsym{→}
    \hyperref[ch-21]{第 21 章}；600-604 \dspsym{→} \hyperref[ch-32]{第
    32 章}
  \item
    notch filter：陷波滤波器 --- 519-520 \dspsym{→} \hyperref[ch-28]{第
    28 章}（疑为旧版页码，按内容应在 565-566 \dspsym{→}
    \hyperref[ch-30]{第 30 章}）；592-597 \dspsym{→} \hyperref[ch-32]{第
    32 章}
  \item
    overshoot, step response：过冲，阶跃响应 --- 54-55 \dspsym{→}
    \hyperref[ch-03]{第 3 章}
  \item
    pulse response：脉冲响应 --- 55 \dspsym{→} \hyperref[ch-03]{第 3 章}
  \item
    reconstruction filter for DAC：DAC 的重建滤波器 --- 44-49 \dspsym{→}
    \hyperref[ch-03]{第 3 章}；480 \dspsym{→} \hyperref[ch-26]{第 26 章}
  \item
    ringing, step response：振铃，阶跃响应 --- 54-55 \dspsym{→}
    \hyperref[ch-03]{第 3 章}
  \item
    roll-off：滚降 --- 52-53 \dspsym{→} \hyperref[ch-03]{第 3 章}
  \item
    Sallen-key circuit：Sallen-Key 电路 --- 49-50 \dspsym{→}
    \hyperref[ch-03]{第 3 章}；600-601 \dspsym{→} \hyperref[ch-32]{第 32
    章}
  \item
    smoothing：平滑 --- 322 \dspsym{→} \hyperref[ch-19]{第 19 章}
  \item
    stability of analog filters：模拟滤波器的稳定性 --- 541 \dspsym{→}
    \hyperref[ch-29]{第 29 章}（疑为旧版页码，按内容应在 589 \dspsym{→}
    \hyperref[ch-32]{第 32 章}）；600-601 \dspsym{→} \hyperref[ch-32]{第
    32 章}
  \item
    step response：阶跃响应 --- 54-55 \dspsym{→} \hyperref[ch-03]{第 3
    章}；322-323 \dspsym{→} \hyperref[ch-19]{第 19 章}
  \item
    switch-capacitor：开关电容 --- 51-52 \dspsym{→} \hyperref[ch-03]{第
    3 章}
  \end{itemize}
\item
  \textbf{Filters, digital} 滤波器，数字

  \begin{itemize}
  \tightlist
  \item
    另见 \textbf{Filters, analog}（滤波器，模拟）；\textbf{Filters,
    recursive}（滤波器，递归）
  \item
    band-pass：带通 --- 177-179 \dspsym{→} \hyperref[ch-09]{第 9
    章}；268-274 \dspsym{→} \hyperref[ch-14]{第 14 章}；293 \dspsym{→}
    \hyperref[ch-16]{第 16 章}
  \item
    band-reject：带阻 --- 268-274 \dspsym{→} \hyperref[ch-14]{第 14
    章}；291, 293 \dspsym{→} \hyperref[ch-16]{第 16 章}
  \item
    custom response：自定义响应 --- 297-310 \dspsym{→}
    \hyperref[ch-17]{第 17 章}
  \item
    cutoff frequency (defined)：截止频率（定义） --- 268 \dspsym{→}
    \hyperref[ch-14]{第 14 章}
  \item
    edge enhancement：边缘增强 --- 400-401 \dspsym{→}
    \hyperref[ch-24]{第 24 章}
  \item
    even order filter：偶数阶滤波器 --- 603 \dspsym{→}
    \hyperref[ch-32]{第 32 章}
  \item
    finite impulse response (FIR)：有限冲激响应（FIR） --- 263
    \dspsym{→} \hyperref[ch-14]{第 14 章}；319 \dspsym{→}
    \hyperref[ch-19]{第 19 章}
  \item
    FIR vs.~IIR：FIR 与 IIR 对比 --- 346-350 \dspsym{→}
    \hyperref[ch-21]{第 21 章}
  \item
    frequency domain parameters：频域参数 --- 268-270 \dspsym{→}
    \hyperref[ch-14]{第 14 章}
  \item
    frequency response：频率响应 --- 176-179 \dspsym{→}
    \hyperref[ch-09]{第 9 章}；268-270 \dspsym{→} \hyperref[ch-14]{第 14
    章}；562 \dspsym{→} \hyperref[ch-30]{第 30 章}
  \item
    high-pass：高通 --- 110 \dspsym{→} \hyperref[ch-06]{第 6 章}；129
    \dspsym{→} \hyperref[ch-07]{第 7 章}；268-274 \dspsym{→}
    \hyperref[ch-14]{第 14 章}
  \item
    low-pass：低通 --- 110 \dspsym{→} \hyperref[ch-06]{第 6 章}；128-129
    \dspsym{→} \hyperref[ch-07]{第 7 章}；268-270 \dspsym{→}
    \hyperref[ch-14]{第 14 章}；280 \dspsym{→} \hyperref[ch-15]{第 15
    章}；285-288 \dspsym{→} \hyperref[ch-16]{第 16 章}；343-350
    \dspsym{→} \hyperref[ch-21]{第 21 章}；400-401 \dspsym{→}
    \hyperref[ch-24]{第 24 章}
  \item
    matched filter：匹配滤波器 --- 138 \dspsym{→} \hyperref[ch-07]{第 7
    章}；307-310 \dspsym{→} \hyperref[ch-17]{第 17 章}；416-419
    \dspsym{→} \hyperref[ch-24]{第 24 章}
  \item
    moving average filter：移动平均滤波器 --- 277-284 \dspsym{→}
    \hyperref[ch-15]{第 15 章}；307-310 \dspsym{→} \hyperref[ch-17]{第
    17 章}；348-350 \dspsym{→} \hyperref[ch-21]{第 21 章}；406-407
    \dspsym{→} \hyperref[ch-24]{第 24 章}
  \item
    odd order filter：奇数阶滤波器 --- 603 \dspsym{→}
    \hyperref[ch-32]{第 32 章}
  \item
    overshoot, step response：过冲，阶跃响应 --- 266-267 \dspsym{→}
    \hyperref[ch-14]{第 14 章}
  \item
    passband (defined)：通带（定义） --- 268 \dspsym{→}
    \hyperref[ch-14]{第 14 章}
  \item
    passband ripple (defined)：通带纹波（定义） --- 268 \dspsym{→}
    \hyperref[ch-14]{第 14 章}
  \item
    pulse response：脉冲响应 --- 328-332 \dspsym{→} \hyperref[ch-19]{第
    19 章}
  \item
    roll-off (defined)：滚降（定义） --- 268-269 \dspsym{→}
    \hyperref[ch-14]{第 14 章}
  \item
    smoothing：平滑 --- 110 \dspsym{→} \hyperref[ch-06]{第 6
    章}；128-129 \dspsym{→} \hyperref[ch-07]{第 7 章}；280-282
    \dspsym{→} \hyperref[ch-15]{第 15 章}；348-350 \dspsym{→}
    \hyperref[ch-21]{第 21 章}
  \item
    spectral inversion：频谱反转 --- 270-272 \dspsym{→}
    \hyperref[ch-14]{第 14 章}；293 \dspsym{→} \hyperref[ch-16]{第 16
    章}
  \item
    spectral reversal：频谱翻转 --- 273-274 \dspsym{→}
    \hyperref[ch-14]{第 14 章}
  \item
    step response：阶跃响应 --- 262-263, 265-267 \dspsym{→}
    \hyperref[ch-14]{第 14 章}；338 \dspsym{→} \hyperref[ch-20]{第 20
    章}
  \item
    stopband (defined)：阻带（定义） --- 268 \dspsym{→}
    \hyperref[ch-14]{第 14 章}
  \item
    stopband attenuation：阻带衰减 --- 268-269 \dspsym{→}
    \hyperref[ch-14]{第 14 章}；293-296 \dspsym{→} \hyperref[ch-16]{第
    16 章}
  \item
    time domain parameters：时域参数 --- 266-268 \dspsym{→}
    \hyperref[ch-14]{第 14 章}
  \item
    transition band (defined)：过渡带（定义） --- 268 \dspsym{→}
    \hyperref[ch-14]{第 14 章}
  \item
    Wiener filter：维纳滤波器 --- 308-310 \dspsym{→} \hyperref[ch-17]{第
    17 章}；368-370 \dspsym{→} \hyperref[ch-22]{第 22 章}
  \item
    windowed-sinc filter：加窗 sinc 滤波器 --- 216-217 \dspsym{→}
    \hyperref[ch-11]{第 11 章}；285-298 \dspsym{→} \hyperref[ch-16]{第
    16 章}、\hyperref[ch-17]{第 17 章}；346-348 \dspsym{→}
    \hyperref[ch-21]{第 21 章}
  \end{itemize}
\item
  \textbf{Filters, recursive} 滤波器，递归

  \begin{itemize}
  \tightlist
  \item
    另见 \textbf{Filters, analog}（滤波器，模拟）；\textbf{Filters,
    digital}（滤波器，数字）
  \item
    band-pass filters：带通滤波器 --- 326-327 \dspsym{→}
    \hyperref[ch-19]{第 19 章}
  \item
    band-reject filters：带阻滤波器 --- 326-327 \dspsym{→}
    \hyperref[ch-19]{第 19 章}；610-616 \dspsym{→} \hyperref[ch-33]{第
    33 章}
  \item
    bidirectional recursive filtering：双向递归滤波 --- 330-332
    \dspsym{→} \hyperref[ch-19]{第 19 章}
  \item
    Butterworth：巴特沃斯 --- 333-342 \dspsym{→} \hyperref[ch-20]{第 20
    章}；601-603 \dspsym{→} \hyperref[ch-32]{第 32 章}；623-630
    \dspsym{→} \hyperref[ch-33]{第 33 章}
  \item
    cascade stages, combining：合并级联级 --- 616-618 \dspsym{→}
    \hyperref[ch-33]{第 33 章}
  \item
    Chebyshev：切比雪夫 --- 333-342 \dspsym{→} \hyperref[ch-20]{第 20
    章}；346-348 \dspsym{→} \hyperref[ch-21]{第 21 章}；623-630
    \dspsym{→} \hyperref[ch-33]{第 33 章}
  \item
    converting pole and zeros to recursion
    coefficients：把极点和零点转换成递归系数 --- 610-615 \dspsym{→}
    \hyperref[ch-33]{第 33 章}
  \item
    custom response：自定义响应（迭代设计） --- 476-480 \dspsym{→}
    \hyperref[ch-26]{第 26 章}
  \item
    elliptic：椭圆 --- 334 \dspsym{→} \hyperref[ch-20]{第 20
    章}；602-603 \dspsym{→} \hyperref[ch-32]{第 32 章}
  \item
    FIR compared with IIR：FIR 与 IIR 比较 --- 346-348 \dspsym{→}
    \hyperref[ch-21]{第 21 章}
  \item
    gain changes：增益调整 --- 621-623 \dspsym{→} \hyperref[ch-33]{第 33
    章}
  \item
    high-pass：高通 --- 322-326 \dspsym{→} \hyperref[ch-19]{第 19
    章}；334-338 \dspsym{→} \hyperref[ch-20]{第 20 章}
  \item
    infinite impulse response (IIR)：无限冲激响应（IIR） --- 263
    \dspsym{→} \hyperref[ch-14]{第 14 章}；284 \dspsym{→}
    \hyperref[ch-15]{第 15 章}；319-320 \dspsym{→} \hyperref[ch-19]{第
    19 章}
  \item
    low-pass：低通 --- 322-326 \dspsym{→} \hyperref[ch-19]{第 19
    章}；334-338 \dspsym{→} \hyperref[ch-20]{第 20 章}；346-350
    \dspsym{→} \hyperref[ch-21]{第 21 章}
  \item
    low-pass to low-pass transform：低通到低通变换 --- 628-629
    \dspsym{→} \hyperref[ch-33]{第 33 章}
  \item
    low-pass to high-pass transform：低通到高通变换 --- 629 \dspsym{→}
    \hyperref[ch-33]{第 33 章}
  \item
    moving average filter, recursive：移动平均滤波器，递归实现 ---
    282-284 \dspsym{→} \hyperref[ch-15]{第 15 章}
  \item
    narrow-band filters：窄带滤波器 --- 326-327 \dspsym{→}
    \hyperref[ch-19]{第 19 章}
  \item
    notch filters：陷波滤波器 --- 326-327 \dspsym{→} \hyperref[ch-19]{第
    19 章}；612-614 \dspsym{→} \hyperref[ch-33]{第 33 章}
  \item
    overshoot, step response：过冲，阶跃响应 --- 338 \dspsym{→}
    \hyperref[ch-20]{第 20 章}
  \item
    parallel stages, combining：合并并联级 --- 616-619 \dspsym{→}
    \hyperref[ch-33]{第 33 章}
  \item
    pulse response：脉冲响应 --- 328-332 \dspsym{→} \hyperref[ch-19]{第
    19 章}
  \item
    reconstruction filter for DAC：DAC 的重建滤波器 --- 44-55 \dspsym{→}
    \hyperref[ch-03]{第 3 章}；480 \dspsym{→} \hyperref[ch-26]{第 26 章}
  \item
    ringing, step response：振铃，阶跃响应 --- 338 \dspsym{→}
    \hyperref[ch-20]{第 20 章}
  \item
    single pole recursive filters：单极点递归滤波器 --- 322-326
    \dspsym{→} \hyperref[ch-19]{第 19 章}；348-350 \dspsym{→}
    \hyperref[ch-21]{第 21 章}；406-407 \dspsym{→} \hyperref[ch-24]{第
    24 章}
  \item
    smoothing：平滑 --- 322-326 \dspsym{→} \hyperref[ch-19]{第 19
    章}；348-350 \dspsym{→} \hyperref[ch-21]{第 21 章}
  \item
    spectral inversion：频谱反转 --- 619-621 \dspsym{→}
    \hyperref[ch-33]{第 33 章}
  \item
    stability of recursive filters：递归滤波器的稳定性 --- 339
    \dspsym{→} \hyperref[ch-20]{第 20 章}；599, 598-599 \dspsym{→}
    \hyperref[ch-32]{第 32 章}；609 \dspsym{→} \hyperref[ch-33]{第 33
    章}
  \item
    step response：阶跃响应 --- 322-323 \dspsym{→} \hyperref[ch-19]{第
    19 章}；338 \dspsym{→} \hyperref[ch-20]{第 20 章}
  \item
    transfer function of：其传递函数 --- 610-616 \dspsym{→}
    \hyperref[ch-33]{第 33 章}
  \item
    z-domain representation of：其 z 域表示 --- 610-616 \dspsym{→}
    \hyperref[ch-33]{第 33 章}
  \end{itemize}
\item
  \textbf{Filter kernel (defined)} 滤波器核（定义） --- 262 \dspsym{→}
  \hyperref[ch-14]{第 14 章}
\item
  \textbf{Filtered backprojection} 滤波反投影 --- 446-450 \dspsym{→}
  \hyperref[ch-25]{第 25 章}
\item
  \textbf{Fingerprint identification} 指纹识别 --- 438-442 \dspsym{→}
  \hyperref[ch-25]{第 25 章}
\item
  \textbf{Finite impulse response} 有限冲激响应 --- 见 \textbf{Filters,
  digital}（滤波器，数字）
\item
  \textbf{FIR filters} FIR 滤波器 --- 见 \textbf{Filters,
  digital}（滤波器，数字）
\item
  \textbf{First difference} 一阶差分 --- 110-111 \dspsym{→}
  \hyperref[ch-06]{第 6 章}；125-128 \dspsym{→} \hyperref[ch-07]{第 7
  章}
\item
  \textbf{Fixed point} 定点 --- 68-70 \dspsym{→} \hyperref[ch-04]{第 4
  章}；514 \dspsym{→} \hyperref[ch-28]{第 28 章}；544 \dspsym{→}
  \hyperref[ch-29]{第 29 章}
\item
  \textbf{Flat-top window} 平顶窗 --- 176 \dspsym{→} \hyperref[ch-09]{第
  9 章}
\item
  \textbf{Floating point} 浮点 --- 70-72 \dspsym{→} \hyperref[ch-04]{第
  4 章}；514 \dspsym{→} \hyperref[ch-28]{第 28 章}；544 \dspsym{→}
  \hyperref[ch-29]{第 29 章}
\item
  \textbf{Focusing, eyes and camera} 对焦，眼睛与相机 --- 376-379
  \dspsym{→} \hyperref[ch-23]{第 23 章}
\item
  \textbf{Format frequencies in speech} 语音中的共振峰频率 --- 365
  \dspsym{→} \hyperref[ch-22]{第 22 章}
\item
  \textbf{Fovea of the eye} 眼睛的中央凹 --- 379-380 \dspsym{→}
  \hyperref[ch-23]{第 23 章}
\item
  \textbf{Fourier slice theorem} 傅里叶切片定理 --- 411 \dspsym{→}
  \hyperref[ch-24]{第 24 章}；448-449 \dspsym{→} \hyperref[ch-25]{第 25
  章}
\item
  \textbf{Fourier transform} 傅里叶变换

  \begin{itemize}
  \tightlist
  \item
    另见 \textbf{Discrete Fourier
    transform}（离散傅里叶变换）；\textbf{Fourier transform
    pairs}（傅里叶变换对）
  \item
    analysis equations：分析方程 --- 147, 157-161 \dspsym{→}
    \hyperref[ch-08]{第 8 章}；576-579 \dspsym{→} \hyperref[ch-31]{第 31
    章}
  \item
    circular：循环 --- 见 \textbf{Discrete Fourier
    transform}（离散傅里叶变换）条下
  \item
    data compression use of：在数据压缩中的应用 --- 494-495 \dspsym{→}
    \hyperref[ch-27]{第 27 章}
  \item
    decomposition：分解 --- 见 \textbf{Decomposition}（分解）
  \item
    DFT：DFT --- 见 \textbf{Discrete Fourier
    Transform}（离散傅里叶变换）
  \item
    discrete vs.~periodic relationship：离散与周期的对应关系 --- 222
    \dspsym{→} \hyperref[ch-11]{第 11 章}
  \item
    discrete time Fourier series (DTFT)：离散时间傅里叶级数（DTFT） ---
    144-145 \dspsym{→} \hyperref[ch-08]{第 8 章}；178 \dspsym{→}
    \hyperref[ch-09]{第 9 章}；206-207 \dspsym{→} \hyperref[ch-10]{第 10
    章}；213 \dspsym{→} \hyperref[ch-11]{第 11 章}；530-531 \dspsym{→}
    \hyperref[ch-28]{第 28 章}
  \item
    four types of：四种类型 --- 143-145 \dspsym{→} \hyperref[ch-08]{第 8
    章}；576-579 \dspsym{→} \hyperref[ch-31]{第 31 章}
  \item
    Fourier series：傅里叶级数 --- 144-145 \dspsym{→}
    \hyperref[ch-08]{第 8 章}；252, 255-260 \dspsym{→}
    \hyperref[ch-13]{第 13 章}；578-579 \dspsym{→} \hyperref[ch-31]{第
    31 章}
  \item
    forward transform (defined)：正变换（定义） --- 147 \dspsym{→}
    \hyperref[ch-08]{第 8 章}
  \item
    Fourier transform, continuous：连续傅里叶变换 --- 144-145 \dspsym{→}
    \hyperref[ch-08]{第 8 章}；252-255 \dspsym{→} \hyperref[ch-13]{第 13
    章}；579 \dspsym{→} \hyperref[ch-31]{第 31 章}
  \item
    images, Fourier transform of：图像的傅里叶变换 --- 410-416
    \dspsym{→} \hyperref[ch-24]{第 24 章}
  \item
    inverse transform (defined)：逆变换（定义） --- 147 \dspsym{→}
    \hyperref[ch-08]{第 8 章}
  \item
    Jean Baptiste Joseph Fourier：让·巴蒂斯特·约瑟夫·傅里叶 --- 141
    \dspsym{→} \hyperref[ch-08]{第 8 章}
  \item
    neural networks, can be carried out by：可以用神经网络实现 ---
    464-465 \dspsym{→} \hyperref[ch-26]{第 26 章}
  \item
    periodic nature：周期性 --- 见 \textbf{Discrete Fourier
    transform}（离散傅里叶变换）条下
  \item
    real vs.~complex：实数与复数形式对比 --- 146 \dspsym{→}
    \hyperref[ch-08]{第 8 章}；225-227 \dspsym{→} \hyperref[ch-12]{第 12
    章}；576 \dspsym{→} \hyperref[ch-31]{第 31 章}
  \item
    real and imaginary parts：实部与虚部 --- 148 \dspsym{→}
    \hyperref[ch-08]{第 8 章}
  \item
    scaling：缩放 --- 529-531 \dspsym{→} \hyperref[ch-28]{第 28
    章}（疑为旧版页码，按内容应在 577-579 \dspsym{→} \hyperref[ch-31]{第
    31 章}）
  \item
    synthesis equations：合成方程 --- 152-157 \dspsym{→}
    \hyperref[ch-08]{第 8 章}；577-579 \dspsym{→} \hyperref[ch-31]{第 31
    章}
  \item
    why sinusoids are used：为什么用正弦波 --- 142 \dspsym{→}
    \hyperref[ch-08]{第 8 章}
  \end{itemize}
\item
  \textbf{Fourier Transform pairs} 傅里叶变换对 --- 209-224 \dspsym{→}
  \hyperref[ch-11]{第 11 章}

  \begin{itemize}
  \tightlist
  \item
    chirp signals and systems：啁啾信号与啁啾系统 --- 222-224 \dspsym{→}
    \hyperref[ch-11]{第 11 章}
  \item
    delta function：冲激函数 --- 210-211 \dspsym{→} \hyperref[ch-11]{第
    11 章}
  \item
    distorted sinusoid：失真的正弦波 --- 220-221 \dspsym{→}
    \hyperref[ch-11]{第 11 章}
  \item
    Gaussian：高斯 --- 216-217 \dspsym{→} \hyperref[ch-11]{第 11 章}
  \item
    Gaussian burst：高斯脉冲串 --- 216-217 \dspsym{→}
    \hyperref[ch-11]{第 11 章}
  \item
    rectangular pulse：矩形脉冲 --- 212-213 \dspsym{→}
    \hyperref[ch-11]{第 11 章}
  \item
    shifted impulse：平移的冲激 --- 210-211 \dspsym{→}
    \hyperref[ch-11]{第 11 章}
  \item
    sinc：sinc --- 215-217 \dspsym{→} \hyperref[ch-11]{第 11 章}
  \item
    triangular pulse：三角形脉冲 --- 216-217 \dspsym{→}
    \hyperref[ch-11]{第 11 章}
  \end{itemize}
\item
  \textbf{Fourier reconstruction (CT)} 傅里叶重建（CT） --- 447-450
  \dspsym{→} \hyperref[ch-25]{第 25 章}
\item
  \textbf{Frame grabber} 帧采集卡 --- 385 \dspsym{→} \hyperref[ch-23]{第
  23 章}
\item
  \textbf{Frame, television} 帧，电视 --- 384-385 \dspsym{→}
  \hyperref[ch-23]{第 23 章}
\item
  \textbf{Frequency domain (defined)} 频域（定义） --- 12 \dspsym{→}
  \hyperref[ch-02]{第 2 章}；147 \dspsym{→} \hyperref[ch-08]{第 8 章}
\item
  \textbf{Frequency domain encoding} 频域编码 --- 见
  \textbf{Information}（信息）
\item
  \textbf{Frequency domain multiplexing} 频分复用 --- 206 \dspsym{→}
  \hyperref[ch-10]{第 10 章}
\item
  \textbf{Frequency response} 频率响应 --- 见 \textbf{Filters}（滤波器）
\item
  \textbf{Fricative sound in speech} 语音中的摩擦音 --- 364-365
  \dspsym{→} \hyperref[ch-22]{第 22 章}
\item
  \textbf{Full-width-half-maximum (FWHM)} 半高全宽 --- 424 \dspsym{→}
  \hyperref[ch-25]{第 25 章}
\item
  \textbf{Fully interconnected neural network} 全互连神经网络 --- 460
  \dspsym{→} \hyperref[ch-26]{第 26 章}
\item
  \textbf{Fundamental frequency} 基频 --- 见 \textbf{Harmonics}（谐波）
\end{itemize}

\section{G}

\begin{itemize}
\tightlist
\item
  \textbf{Gamma curve} 伽马曲线 --- 389-394 \dspsym{→}
  \hyperref[ch-23]{第 23 章}
\item
  \textbf{Gamma ray detector} 伽马射线探测器 --- 301-305 \dspsym{→}
  \hyperref[ch-17]{第 17 章}
\item
  \textbf{Gauss distribution} 高斯分布 --- 见 \textbf{Gaussian}（高斯）
\item
  \textbf{Gaussian} 高斯

  \begin{itemize}
  \tightlist
  \item
    另见 \textbf{Central limit theorem}（中心极限定理）
  \item
    equation for：公式 --- 26-31 \dspsym{→} \hyperref[ch-02]{第 2 章}
  \item
    Fourier transform of：其傅里叶变换 --- 216-217 \dspsym{→}
    \hyperref[ch-11]{第 11 章}；423-425, 427-428 \dspsym{→}
    \hyperref[ch-25]{第 25 章}
  \item
    as filter kernel：作为滤波器核 --- 281-283 \dspsym{→}
    \hyperref[ch-15]{第 15 章}；400-401 \dspsym{→} \hyperref[ch-24]{第
    24 章}；423-425, 427-428, 430-432 \dspsym{→} \hyperref[ch-25]{第 25
    章}
  \item
    noise：高斯噪声 --- 29-32 \dspsym{→} \hyperref[ch-02]{第 2 章}
  \item
    separability of：其可分离性 --- 404-406 \dspsym{→}
    \hyperref[ch-24]{第 24 章}
  \end{itemize}
\item
  \textbf{Geophysics} 地球物理学 --- 见 \textbf{Seismology}（地震学）
\item
  \textbf{Gibbs effect} 吉布斯效应 --- 142 \dspsym{→}
  \hyperref[ch-08]{第 8 章}；203 \dspsym{→} \hyperref[ch-10]{第 10
  章}；218-219 \dspsym{→} \hyperref[ch-11]{第 11 章}；286-287 \dspsym{→}
  \hyperref[ch-16]{第 16 章}
\item
  \textbf{GIF image, data compression} GIF 图像，数据压缩 --- 488
  \dspsym{→} \hyperref[ch-27]{第 27 章}
\item
  \textbf{Grayscale, image} 灰度，图像 --- 373, 387 \dspsym{→}
  \hyperref[ch-23]{第 23 章}
\item
  \textbf{Grayscale stretch} 灰度拉伸 --- 390-391 \dspsym{→}
  \hyperref[ch-23]{第 23 章}
\item
  \textbf{Grayscale transforms} 灰度变换 --- 390-394 \dspsym{→}
  \hyperref[ch-23]{第 23 章}；433 \dspsym{→} \hyperref[ch-25]{第 25 章}
\end{itemize}

\section{H}

\begin{itemize}
\tightlist
\item
  \textbf{Halftone image} 半色调图像 --- 387 \dspsym{→}
  \hyperref[ch-23]{第 23 章}；433 \dspsym{→} \hyperref[ch-25]{第 25 章}
\item
  \textbf{Harmonics} 谐波 --- 172-173 \dspsym{→} \hyperref[ch-09]{第 9
  章}；220-222 \dspsym{→} \hyperref[ch-11]{第 11 章}；255 \dspsym{→}
  \hyperref[ch-13]{第 13 章}；355-358 \dspsym{→} \hyperref[ch-22]{第 22
  章}
\item
  \textbf{Harvard Architecture} 哈佛结构 --- 84 \dspsym{→}
  \hyperref[ch-04]{第 4 章}；509 \dspsym{→} \hyperref[ch-28]{第 28 章}
\item
  \textbf{Hearing} 听觉 --- 351-355 \dspsym{→} \hyperref[ch-22]{第 22
  章}
\item
  \textbf{High fidelity audio} 高保真音频 --- 358-362 \dspsym{→}
  \hyperref[ch-22]{第 22 章}
\item
  \textbf{High-pass filters} 高通滤波器 --- 见
  \textbf{Filters}（滤波器）
\item
  \textbf{High speed convolution} 高速卷积 --- 见 \textbf{FFT
  convolution}（FFT 卷积）
\item
  \textbf{Hilbert transformer} 希尔伯特变换器 --- 621 \dspsym{→}
  \hyperref[ch-33]{第 33 章}
\item
  \textbf{Histogram, calculating} 直方图，计算 --- 19-26 \dspsym{→}
  \hyperref[ch-02]{第 2 章}
\item
  \textbf{Histogram equalization} 直方图均衡化 --- 393-394 \dspsym{→}
  \hyperref[ch-23]{第 23 章}
\item
  \textbf{History of DSP} DSP 的历史 --- 1-3 \dspsym{→}
  \hyperref[ch-01]{第 1 章}
\item
  \textbf{Homogeneity} 齐次性 --- 89-90 \dspsym{→} \hyperref[ch-05]{第 5
  章}；108 \dspsym{→} \hyperref[ch-06]{第 6 章}；185-186 \dspsym{→}
  \hyperref[ch-10]{第 10 章}
\item
  \textbf{Homomorphic processing} 同态处理 --- 370-372 \dspsym{→}
  \hyperref[ch-22]{第 22 章}；408-409 \dspsym{→} \hyperref[ch-24]{第 24
  章}
\item
  \textbf{Hyperspace} 超空间 --- 457, 463-464 \dspsym{→}
  \hyperref[ch-26]{第 26 章}
\item
  \textbf{Huffman encoding} 霍夫曼编码 --- 484-486, 500 \dspsym{→}
  \hyperref[ch-27]{第 27 章}
\end{itemize}

\section{I}

\begin{itemize}
\tightlist
\item
  \textbf{IIR filters} IIR 滤波器 --- 见 \textbf{Filters,
  recursive}（滤波器，递归）
\item
  \textbf{Illumination flattening of images} 图像的照明平坦化 ---
  407-410 \dspsym{→} \hyperref[ch-24]{第 24 章}
\item
  \textbf{Immersion} 浸没 --- 见 \textbf{Convolution}（卷积）条下
\item
  \textbf{Impedance, electrical} 阻抗，电 --- 563-566 \dspsym{→}
  \hyperref[ch-30]{第 30 章}
\item
  \textbf{Impulse} 冲激 --- 100 \dspsym{→} \hyperref[ch-05]{第 5
  章}；107-108 \dspsym{→} \hyperref[ch-06]{第 6 章}
\item
  \textbf{Impulse decomposition} 冲激分解 --- 见
  \textbf{decomposition}（分解）条下
\item
  \textbf{Impulse response} 冲激响应

  \begin{itemize}
  \tightlist
  \item
    另见 \textbf{Convolution}（卷积）
  \item
    continuous systems：连续系统 --- 244-245 \dspsym{→}
    \hyperref[ch-13]{第 13 章}
  \item
    defined：定义 --- 108-109 \dspsym{→} \hyperref[ch-06]{第 6 章}
  \item
    examples of：示例 --- 110-111 \dspsym{→} \hyperref[ch-06]{第 6
    章}；128-132 \dspsym{→} \hyperref[ch-07]{第 7 章}
  \item
    two-dimensional (image)：二维（图像） --- 397-398 \dspsym{→}
    \hyperref[ch-24]{第 24 章}
  \end{itemize}
\item
  \textbf{Impulse train} 冲激串 --- 42-47 \dspsym{→} \hyperref[ch-03]{第
  3 章}
\item
  \textbf{Independent variable (defined)} 自变量（定义） --- 12
  \dspsym{→} \hyperref[ch-02]{第 2 章}
\item
  \textbf{Infinite impulse response (IIR) filters} 无限冲激响应滤波器
  --- 见 \textbf{Filters, recursive}（滤波器，递归）
\item
  \textbf{Information} 信息

  \begin{itemize}
  \tightlist
  \item
    frequency domain encoded：频域编码 --- 56 \dspsym{→}
    \hyperref[ch-03]{第 3 章}；265-266, 268-270 \dspsym{→}
    \hyperref[ch-14]{第 14 章}
  \item
    spatial domain encoded：空间域编码 --- 373-374 \dspsym{→}
    \hyperref[ch-23]{第 23 章}；424-425 \dspsym{→} \hyperref[ch-25]{第
    25 章}
  \item
    time domain encoded：时域编码 --- 56-57 \dspsym{→}
    \hyperref[ch-03]{第 3 章}；265-267 \dspsym{→} \hyperref[ch-14]{第 14
    章}
  \end{itemize}
\item
  \textbf{In-place computation} 原地计算 --- 233 \dspsym{→}
  \hyperref[ch-12]{第 12 章}
\item
  \textbf{Integral, discrete} 积分，离散 --- 见 \textbf{Running
  sum}（累加和）
\item
  \textbf{Integral of continuous impulse} 连续冲激的积分 --- 244,
  250-251 \dspsym{→} \hyperref[ch-13]{第 13 章}
\item
  \textbf{Integrated profile} 积分剖面 --- 428-429 \dspsym{→}
  \hyperref[ch-25]{第 25 章}
\item
  \textbf{Interlaced decomposition} 交错分解 --- 见
  \textbf{Decomposition}（分解）条下
\item
  \textbf{Interlaced video} 隔行视频 --- 384-385 \dspsym{→}
  \hyperref[ch-23]{第 23 章}
\item
  \textbf{Interpolation} 插值 --- 见 \textbf{Multirate}
\item
  \textbf{Iterative techniques} 迭代技术 --- 444-445 \dspsym{→}
  \hyperref[ch-25]{第 25 章}；465-473 \dspsym{→} \hyperref[ch-26]{第 26
  章}
\item
  \textbf{Iterative least squares technique (ILST)} 迭代最小二乘技术 ---
  444 \dspsym{→} \hyperref[ch-25]{第 25 章}
\end{itemize}

\section{J}

\begin{itemize}
\tightlist
\item
  \textbf{JPEG image compression} JPEG 图像压缩 --- 494-502 \dspsym{→}
  \hyperref[ch-27]{第 27 章}
\end{itemize}

\section{K}

\begin{itemize}
\tightlist
\item
  \textbf{Karhunen-Loeve transform} 卡南-洛伊夫变换 --- 496 \dspsym{→}
  \hyperref[ch-27]{第 27 章}
\item
  \textbf{Kernel, filter} 核，滤波器 --- 见 \textbf{Impulse
  response}（冲激响应）
\end{itemize}

\section{L}

\begin{itemize}
\tightlist
\item
  \textbf{Laplace transform} 拉普拉斯变换 --- 334 \dspsym{→}
  \hyperref[ch-20]{第 20 章}；581-604 \dspsym{→} \hyperref[ch-32]{第 32
  章}
\item
  \textbf{Layers, neural network} 层，神经网络 --- 459-460 \dspsym{→}
  \hyperref[ch-26]{第 26 章}
\item
  \textbf{Learning algorithm, neural network} 学习算法，神经网络 --- 463
  \dspsym{→} \hyperref[ch-26]{第 26 章}
\item
  \textbf{Least significant bit (LSB) (defined)} 最低有效位（定义） ---
  36 \dspsym{→} \hyperref[ch-03]{第 3 章}
\item
  \textbf{Lens, camera and eye} 透镜，相机与眼睛 --- 376-379 \dspsym{→}
  \hyperref[ch-23]{第 23 章}
\item
  \textbf{Limiting resolution, images} 极限分辨率，图像 --- 426
  \dspsym{→} \hyperref[ch-25]{第 25 章}
\item
  \textbf{Line pair} 线对 --- 426 \dspsym{→} \hyperref[ch-25]{第 25 章}
\item
  \textbf{Line pair gauge} 线对卡 --- 425-426 \dspsym{→}
  \hyperref[ch-25]{第 25 章}
\item
  \textbf{Line scanning image acquisition} 线扫描图像采集 --- 386-387
  \dspsym{→} \hyperref[ch-23]{第 23 章}
\item
  \textbf{Line spread function (LSF)} 线扩散函数 --- 426-430 \dspsym{→}
  \hyperref[ch-25]{第 25 章}
\item
  \textbf{Linear phase} 线性相位 --- 见 \textbf{Phase}（相位）条下
\item
  \textbf{Linear predictive coding (LPC)} 线性预测编码 --- 359, 366
  \dspsym{→} \hyperref[ch-22]{第 22 章}
\item
  \textbf{Linear systems} 线性系统 --- 见 \textbf{Linearity}（线性）
\item
  \textbf{Linearity} 线性

  \begin{itemize}
  \tightlist
  \item
    alternatives：替代方案 --- 104-106 \dspsym{→} \hyperref[ch-05]{第 5
    章}
  \item
    commutative property：交换律 --- 96 \dspsym{→} \hyperref[ch-05]{第 5
    章}
  \item
    decomposition：分解 --- 见 \textbf{Decomposition}（分解）
  \item
    examples of linear and nonlinear systems：线性与非线性系统的例子 ---
    94-95 \dspsym{→} \hyperref[ch-05]{第 5 章}
  \item
    Fourier transform, of the：傅里叶变换的线性 --- 185-188 \dspsym{→}
    \hyperref[ch-10]{第 10 章}
  \item
    memoryless systems：无记忆系统 --- 93 \dspsym{→} \hyperref[ch-05]{第
    5 章}
  \item
    multiplication：乘法 --- 97 \dspsym{→} \hyperref[ch-05]{第 5 章}
  \item
    noise, adding：加入噪声 --- 97 \dspsym{→} \hyperref[ch-05]{第 5 章}
  \item
    requirements for：所需条件 --- 89 \dspsym{→} \hyperref[ch-05]{第 5
    章}
  \item
    sinusoidal fidelity：正弦保真性 --- 92-94 \dspsym{→}
    \hyperref[ch-05]{第 5 章}；142 \dspsym{→} \hyperref[ch-08]{第 8 章}
  \item
    static linearity：静态线性 --- 92-93 \dspsym{→} \hyperref[ch-05]{第
    5 章}
  \item
    superposition：叠加 --- 98-100 \dspsym{→} \hyperref[ch-05]{第 5 章}
  \item
    synthesis：合成 --- 98-99 \dspsym{→} \hyperref[ch-05]{第 5 章}
  \end{itemize}
\item
  \textbf{Logarithmic scale} 对数刻度 --- 见 \textbf{Decibels}（分贝）
\item
  \textbf{Long integer} 长整型 --- 72 \dspsym{→} \hyperref[ch-04]{第 4
  章}
\item
  \textbf{Lossless data compression} 无损数据压缩 --- 481-494 \dspsym{→}
  \hyperref[ch-27]{第 27 章}
\item
  \textbf{Lossy data compression} 有损数据压缩 --- 481-482, 494-502
  \dspsym{→} \hyperref[ch-27]{第 27 章}
\item
  \textbf{Low-pass filters} 低通滤波器 --- 见
  \textbf{Filters}（滤波器）条下
\item
  \textbf{Luminance signal, television} 亮度信号，电视 --- 386
  \dspsym{→} \hyperref[ch-23]{第 23 章}
\item
  \textbf{LZW encoding, data compression} LZW 编码，数据压缩 --- 488-494
  \dspsym{→} \hyperref[ch-27]{第 27 章}
\end{itemize}

\section{M}

\begin{itemize}
\tightlist
\item
  \textbf{Magnetic resonance imaging (MRI)} 磁共振成像 --- 9-10
  \dspsym{→} \hyperref[ch-01]{第 1 章}；450 \dspsym{→}
  \hyperref[ch-25]{第 25 章}
\item
  \textbf{Magnitude} 幅度 --- 见 \textbf{Polar notation}（极坐标表示法）
\item
  \textbf{Matched filter} 匹配滤波器 --- 见 \textbf{Filters,
  digital}（滤波器，数字）条下
\item
  \textbf{Math coprocessor} 数学协处理器 --- 81 \dspsym{→}
  \hyperref[ch-04]{第 4 章}
\item
  \textbf{Mean} 均值 --- 13-17, 20-22 \dspsym{→} \hyperref[ch-02]{第 2
  章}；434-436 \dspsym{→} \hyperref[ch-25]{第 25 章}
\item
  \textbf{Medical imaging} 医学成像 --- 2 \dspsym{→} \hyperref[ch-01]{第
  1 章}，另见 \textbf{X-ray imaging}（X 射线成像）；\textbf{Computed
  Tomography}（计算机断层成像）
\item
  \textbf{Memory cache} 高速缓存 --- 81-82 \dspsym{→}
  \hyperref[ch-04]{第 4 章}
\item
  \textbf{Memoryless system} 无记忆系统 --- 93 \dspsym{→}
  \hyperref[ch-05]{第 5 章}
\item
  \textbf{MFLOPS} MFLOPS（每秒百万次浮点运算） --- 526 \dspsym{→}
  \hyperref[ch-28]{第 28 章}
\item
  \textbf{Microphonics} 微音效应 --- 172 \dspsym{→} \hyperref[ch-09]{第
  9 章}
\item
  \textbf{MIPS} MIPS（每秒百万条指令） --- 526 \dspsym{→}
  \hyperref[ch-28]{第 28 章}
\item
  \textbf{Mix down, music} 缩混，音乐 --- 5 \dspsym{→}
  \hyperref[ch-01]{第 1 章}；362 \dspsym{→} \hyperref[ch-22]{第 22 章}
\item
  \textbf{Modulation transfer function} 调制传递函数 --- 425-432
  \dspsym{→} \hyperref[ch-25]{第 25 章}
\item
  \textbf{Morphing} 变形动画 --- 394 \dspsym{→} \hyperref[ch-23]{第 23
  章}
\item
  \textbf{Morphological processing} 形态学处理 --- 436-442 \dspsym{→}
  \hyperref[ch-25]{第 25 章}
\item
  \textbf{Moving average filter} 移动平均滤波器 --- 见 \textbf{Filters,
  digital}（滤波器，数字）条下
\item
  \textbf{MPEG} MPEG --- 501-502 \dspsym{→} \hyperref[ch-27]{第 27 章}
\item
  \textbf{MTF} MTF --- 见 \textbf{Modulation transfer
  function}（调制传递函数）
\item
  \textbf{Multiplexing, telephone} 复用，电话 --- 4 \dspsym{→}
  \hyperref[ch-01]{第 1 章}
\item
  \textbf{Multiplication} 乘法

  \begin{itemize}
  \tightlist
  \item
    amplitude modulation：幅度调制 --- 见 \textbf{Amplitude
    modulation}（幅度调制（调幅））
  \item
    frequency domain signals：频域信号 --- 180-181 \dspsym{→}
    \hyperref[ch-09]{第 9 章}；312-316 \dspsym{→} \hyperref[ch-18]{第 18
    章}
  \item
    image formation model：图像形成模型 --- 378 \dspsym{→}
    \hyperref[ch-23]{第 23 章}；407-410 \dspsym{→} \hyperref[ch-24]{第
    24 章}
  \item
    of time domain signals：时域信号的乘法 --- 97 \dspsym{→}
    \hyperref[ch-05]{第 5 章}
  \end{itemize}
\item
  \textbf{Multiprocessing} 多处理 --- 529 \dspsym{→} \hyperref[ch-28]{第
  28 章}
\item
  \textbf{Multirate techniques} 多速率技术

  \begin{itemize}
  \tightlist
  \item
    compact disc DAC：CD 的 DAC --- 361 \dspsym{→} \hyperref[ch-22]{第
    22 章}
  \item
    data conversion：数据转换 --- 58-66 \dspsym{→} \hyperref[ch-03]{第 3
    章}
  \item
    decimation：抽取 --- 60 \dspsym{→} \hyperref[ch-03]{第 3
    章}；202-203 \dspsym{→} \hyperref[ch-10]{第 10 章}
  \item
    interpolation：插值 --- 60 \dspsym{→} \hyperref[ch-03]{第 3
    章}；202-203 \dspsym{→} \hyperref[ch-10]{第 10 章}；361 \dspsym{→}
    \hyperref[ch-22]{第 22 章}
  \item
    single bit ADC and DAC：单比特 ADC 与 DAC --- 60-66 \dspsym{→}
    \hyperref[ch-03]{第 3 章}
  \end{itemize}
\item
  \textbf{Mu law companding} μ 律压扩 --- 362-364 \dspsym{→}
  \hyperref[ch-22]{第 22 章}
\item
  \textbf{Music} 音乐 --- 见 \textbf{Audio processing}（音频处理）
\end{itemize}

\section{N}

\begin{itemize}
\tightlist
\item
  \textbf{Natural frequency} 自然频率 --- 149, 164 \dspsym{→}
  \hyperref[ch-08]{第 8 章}；253 \dspsym{→} \hyperref[ch-13]{第 13 章}
\item
  \textbf{Nearest neighbor rounding} 最近邻取整 --- 396 \dspsym{→}
  \hyperref[ch-23]{第 23 章}
\item
  \textbf{Negative frequencies} 负频率 --- 见 \textbf{Discrete Fourier
  transform}（离散傅里叶变换）条下
\item
  \textbf{Neural networks} 神经网络 --- 368 \dspsym{→}
  \hyperref[ch-22]{第 22 章}；451-480 \dspsym{→} \hyperref[ch-26]{第 26
  章}
\item
  \textbf{Night vision systems} 夜视系统 --- 392 \dspsym{→}
  \hyperref[ch-23]{第 23 章}；424, 436 \dspsym{→} \hyperref[ch-25]{第 25
  章}
\item
  \textbf{Nodes, neural networks} 节点，神经网络 --- 459-462 \dspsym{→}
  \hyperref[ch-26]{第 26 章}
\item
  \textbf{Nonlinear phase} 非线性相位 --- 见 \textbf{Phase}（相位）条下
\item
  \textbf{Noise} 噪声

  \begin{itemize}
  \tightlist
  \item
    1/f noise：1/f 噪声 --- 172-173 \dspsym{→} \hyperref[ch-09]{第 9 章}
  \item
    ADC：ADC --- 见 \textbf{Quantization error}（量化误差）
  \item
    data compression, lossy：有损数据压缩 --- 481-482, 494-502
    \dspsym{→} \hyperref[ch-27]{第 27 章}
  \item
    deconvolution, how noise limits：噪声如何限制反卷积 --- 304
    \dspsym{→} \hyperref[ch-17]{第 17 章}
  \item
    digital generation：数字生成 --- 29-32 \dspsym{→}
    \hyperref[ch-02]{第 2 章}
  \item
    image noise：图像噪声 --- 434-436 \dspsym{→} \hyperref[ch-25]{第 25
    章}
  \item
    in math calculation：数学计算中的 --- 见 \textbf{Round-off
    error}（舍入误差）
  \item
    linearity of added：加性噪声的线性 --- 97 \dspsym{→}
    \hyperref[ch-05]{第 5 章}
  \item
    Poisson noise：泊松噪声 --- 434-436 \dspsym{→} \hyperref[ch-25]{第
    25 章}
  \item
    speech, wideband noise reduction：语音的宽带降噪 --- 368-370
    \dspsym{→} \hyperref[ch-22]{第 22 章}
  \item
    statistical noise：统计噪声 --- 18-19 \dspsym{→} \hyperref[ch-02]{第
    2 章}
  \item
    step response sharpness vs.~noise：阶跃响应陡度与噪声的权衡 ---
    278-279 \dspsym{→} \hyperref[ch-15]{第 15 章}
  \item
    white noise：白噪声 --- 172-173 \dspsym{→} \hyperref[ch-09]{第 9
    章}；307 \dspsym{→} \hyperref[ch-17]{第 17 章}
  \end{itemize}
\item
  \textbf{Normal distribution} 正态分布 --- 见 \textbf{Gaussian}（高斯）
\item
  \textbf{NTSC television} NTSC 电视 --- 386 \dspsym{→}
  \hyperref[ch-23]{第 23 章}
\item
  \textbf{Nyquist rate (frequency)} 奈奎斯特速率（频率） --- 41-42
  \dspsym{→} \hyperref[ch-03]{第 3 章}
\item
  \textbf{Nuclear magnetic resonance imaging (NMR)} 核磁共振成像 --- 450
  \dspsym{→} \hyperref[ch-25]{第 25 章}
\end{itemize}

\section{O}

\begin{itemize}
\tightlist
\item
  \textbf{Octave} 八度 --- 357-358 \dspsym{→} \hyperref[ch-22]{第 22 章}
\item
  \textbf{Odd field in video} 视频中的奇场 --- 384-385 \dspsym{→}
  \hyperref[ch-23]{第 23 章}
\item
  \textbf{Odd symmetry} 奇对称 --- 102 \dspsym{→} \hyperref[ch-05]{第 5
  章}；196 \dspsym{→} \hyperref[ch-10]{第 10 章}；240-241 \dspsym{→}
  \hyperref[ch-12]{第 12 章}
\item
  \textbf{Off-line processing} 离线处理 --- 506 \dspsym{→}
  \hyperref[ch-28]{第 28 章}
\item
  \textbf{Offset binary} 偏移二进制 --- 68-69 \dspsym{→}
  \hyperref[ch-04]{第 4 章}
\item
  \textbf{Oil \& mineral exploration} 石油与矿产勘探 --- 见
  \textbf{Seismology}（地震学）
\item
  \textbf{Opening operation} 开运算 --- 437 \dspsym{→}
  \hyperref[ch-25]{第 25 章}
\item
  \textbf{Optimal filters} 最优滤波器 --- 307-310 \dspsym{→}
  \hyperref[ch-17]{第 17 章}；465 \dspsym{→} \hyperref[ch-26]{第 26 章}
\item
  \textbf{Orthognal basis functions} 正交基函数 --- 158-159 \dspsym{→}
  \hyperref[ch-08]{第 8 章}
\item
  \textbf{Output look-up table (image display)} 输出查找表（图像显示）
  --- 387-389 \dspsym{→} \hyperref[ch-23]{第 23 章}
\item
  \textbf{Output transform (image display)} 输出变换（图像显示） ---
  389-391 \dspsym{→} \hyperref[ch-23]{第 23 章}
\item
  \textbf{Overshoot} 过冲 --- 见 \textbf{Filters}（滤波器）条下，见
  \textbf{Gibbs effect}（吉布斯效应）
\end{itemize}

\section{P}

\begin{itemize}
\tightlist
\item
  \textbf{Packbits, data compression} PackBits，数据压缩 --- 483
  \dspsym{→} \hyperref[ch-27]{第 27 章}
\item
  \textbf{PAL television} PAL 电视 --- 386 \dspsym{→}
  \hyperref[ch-23]{第 23 章}
\item
  \textbf{Parallel processing} 并行处理 --- 529 \dspsym{→}
  \hyperref[ch-28]{第 28 章}
\item
  \textbf{Parallel stages with added outputs} 输出相加的并联级 --- 134
  \dspsym{→} \hyperref[ch-07]{第 7 章}；616-619 \dspsym{→}
  \hyperref[ch-33]{第 33 章}
\item
  \textbf{Parameter space} 参数空间 --- 457 \dspsym{→}
  \hyperref[ch-26]{第 26 章}
\item
  \textbf{Parentheses, used to denote continuous signals}
  圆括号，表示连续信号 --- 87 \dspsym{→} \hyperref[ch-05]{第 5 章}
\item
  \textbf{Parseval's relation} 帕塞瓦尔关系 --- 208 \dspsym{→}
  \hyperref[ch-10]{第 10 章}
\item
  \textbf{Passband} 通带 --- 见 \textbf{Filters}（滤波器）条下
\item
  \textbf{Passband ripple} 通带纹波 --- 见
  \textbf{Filters}（滤波器）条下
\item
  \textbf{Periodic nature of the DFT} DFT 的周期性 --- 见
  \textbf{Discrete Fourier transform}（离散傅里叶变换）条下
\item
  \textbf{Phase} 相位

  \begin{itemize}
  \tightlist
  \item
    另见 \textbf{Polar notation}（极坐标表示法）
  \item
    carries edge information：携带边缘信息 --- 191-192 \dspsym{→}
    \hyperref[ch-10]{第 10 章}；222-224 \dspsym{→} \hyperref[ch-11]{第
    11 章}
  \item
    hearing, insensitivity to：听觉对它不敏感 --- 355-356 \dspsym{→}
    \hyperref[ch-22]{第 22 章}
  \item
    linear phase：线性相位 --- 131-132 \dspsym{→} \hyperref[ch-07]{第 7
    章}；188-191 \dspsym{→} \hyperref[ch-10]{第 10 章}；328-332
    \dspsym{→} \hyperref[ch-19]{第 19 章}
  \item
    nonlinear phase：非线性相位 --- 131-132 \dspsym{→}
    \hyperref[ch-07]{第 7 章}；328-332 \dspsym{→} \hyperref[ch-19]{第 19
    章}
  \item
    nuisances and ambiguities：麻烦与多值性 --- 164-168 \dspsym{→}
    \hyperref[ch-08]{第 8 章}
  \item
    time domain shifting, effect on：时域平移对它的影响 --- 188-191
    \dspsym{→} \hyperref[ch-10]{第 10 章}
  \item
    unwrapping：展开 --- 167 \dspsym{→} \hyperref[ch-08]{第 8
    章}；188-189 \dspsym{→} \hyperref[ch-10]{第 10 章}
  \item
    zero phase：零相位 --- 131-132 \dspsym{→} \hyperref[ch-07]{第 7
    章}；188-191 \dspsym{→} \hyperref[ch-10]{第 10 章}；328-332
    \dspsym{→} \hyperref[ch-19]{第 19 章}
  \end{itemize}
\item
  \textbf{Phase lock loops, linearity of} 锁相环的线性 --- 94 \dspsym{→}
  \hyperref[ch-05]{第 5 章}
\item
  \textbf{Phasor transform} 相量变换 --- 515 \dspsym{→}
  \hyperref[ch-28]{第 28 章}（疑为旧版页码，按内容应在 563 \dspsym{→}
  \hyperref[ch-30]{第 30 章}）
\item
  \textbf{Piano keyboard frequencies} 钢琴键盘的频率 --- 353, 357-358
  \dspsym{→} \hyperref[ch-22]{第 22 章}
\item
  \textbf{Pillbox} 圆柱形（点扩散函数） --- 400-401 \dspsym{→}
  \hyperref[ch-24]{第 24 章}；423-424, 427-428 \dspsym{→}
  \hyperref[ch-25]{第 25 章}
\item
  \textbf{Pipeline} 流水线 --- 84 \dspsym{→} \hyperref[ch-04]{第 4 章}
\item
  \textbf{Pitch} 音高 --- 355-358 \dspsym{→} \hyperref[ch-22]{第 22 章}
\item
  \textbf{Pixel (picture element)} 像素 --- 373 \dspsym{→}
  \hyperref[ch-23]{第 23 章}
\item
  \textbf{Polar notation} 极坐标表示法 --- 161-168 \dspsym{→}
  \hyperref[ch-08]{第 8 章}，另见 \textbf{Phase}（相位）
\item
  \textbf{Polar-to-rectangular conversion} 极坐标到直角坐标的转换 ---
  162 \dspsym{→} \hyperref[ch-08]{第 8 章}
\item
  \textbf{Poles and zeros} 极点与零点 --- 49-50 \dspsym{→}
  \hyperref[ch-03]{第 3 章}；334 \dspsym{→} \hyperref[ch-20]{第 20
  章}；590-604 \dspsym{→} \hyperref[ch-32]{第 32 章}
\item
  \textbf{Point-by-point image acquisition} 逐点图像采集 --- 387
  \dspsym{→} \hyperref[ch-23]{第 23 章}
\item
  \textbf{Point spread function} 点扩散函数 --- 见 \textbf{Impulse
  response}（冲激响应）
\item
  \textbf{Pointer} 指针 --- 507 \dspsym{→} \hyperref[ch-28]{第 28 章}
\item
  \textbf{Poisson statistics} 泊松统计 --- 434-436 \dspsym{→}
  \hyperref[ch-25]{第 25 章}
\item
  \textbf{Positron emission tomography (PET)} 正电子发射断层成像 ---
  449-450 \dspsym{→} \hyperref[ch-25]{第 25 章}
\item
  \textbf{PostScript image, data compression} PostScript 图像，数据压缩
  --- 488 \dspsym{→} \hyperref[ch-27]{第 27 章}
\item
  \textbf{Precision} 精密度 --- 32-34 \dspsym{→} \hyperref[ch-02]{第 2
  章}；68, 72-76 \dspsym{→} \hyperref[ch-04]{第 4 章}
\item
  \textbf{Probability} 概率 --- 17-19 \dspsym{→} \hyperref[ch-02]{第 2
  章}
\item
  \textbf{Probability distribution function (pdf)} 概率分布函数 ---
  19-24 \dspsym{→} \hyperref[ch-02]{第 2 章}；452-455 \dspsym{→}
  \hyperref[ch-26]{第 26 章}
\item
  \textbf{Probability mass function (pmf)} 概率质量函数 --- 18, 19-24
  \dspsym{→} \hyperref[ch-02]{第 2 章}
\end{itemize}

\section{Q}

\begin{itemize}
\tightlist
\item
  \textbf{Quantization error} 量化误差 --- 36-39 \dspsym{→}
  \hyperref[ch-03]{第 3 章}
\item
  \textbf{Quantization levels in images} 图像的量化级 --- 374 \dspsym{→}
  \hyperref[ch-23]{第 23 章}
\item
  \textbf{Quantization table, JPEG} 量化表，JPEG --- 499-500 \dspsym{→}
  \hyperref[ch-27]{第 27 章}
\item
  \textbf{Quantum sink} 量子汇 --- 436 \dspsym{→} \hyperref[ch-25]{第 25
  章}
\end{itemize}

\section{R}

\begin{itemize}
\tightlist
\item
  \textbf{Radar} 雷达 --- 1, 7 \dspsym{→} \hyperref[ch-01]{第 1 章}；88
  \dspsym{→} \hyperref[ch-05]{第 5 章}；137 \dspsym{→}
  \hyperref[ch-07]{第 7 章}；222-224 \dspsym{→} \hyperref[ch-11]{第 11
  章}
\item
  \textbf{Radians} 弧度 --- 见 \textbf{Natural frequency}（自然频率）
\item
  \textbf{Range} 值域 --- 12 \dspsym{→} \hyperref[ch-02]{第 2 章}
\item
  \textbf{Random errors} 随机误差 --- 32-34 \dspsym{→}
  \hyperref[ch-02]{第 2 章}
\item
  \textbf{Random number generator} 随机数生成器 --- 29-32 \dspsym{→}
  \hyperref[ch-02]{第 2 章}；465 \dspsym{→} \hyperref[ch-26]{第 26 章}
\item
  \textbf{Real FFT} 实数 FFT --- 239-242 \dspsym{→} \hyperref[ch-12]{第
  12 章}
\item
  \textbf{Real time processing} 实时处理 --- 311 \dspsym{→}
  \hyperref[ch-18]{第 18 章}；506 \dspsym{→} \hyperref[ch-28]{第 28 章}
\item
  \textbf{Receiver operating characteristic (ROC)} 接收者操作特性 ---
  454-455, 474-475 \dspsym{→} \hyperref[ch-26]{第 26 章}
\item
  \textbf{Reconstruction algorithms, CT} 重建算法，CT --- 444-450
  \dspsym{→} \hyperref[ch-25]{第 25 章}
\item
  \textbf{Reconstruction filter} 重建滤波器 --- 见 \textbf{Filters,
  analog}（滤波器，模拟）条下
\item
  \textbf{Rectangular-to-polar conversion} 直角坐标到极坐标的转换 ---
  162 \dspsym{→} \hyperref[ch-08]{第 8 章}
\item
  \textbf{Recursion equation} 递推方程

  \begin{itemize}
  \tightlist
  \item
    另见 \textbf{Filters, recursive}（滤波器，递归）
  \item
    first difference：一阶差分 --- 125-128 \dspsym{→}
    \hyperref[ch-07]{第 7 章}
  \item
    running sum：累加和 --- 125-128 \dspsym{→} \hyperref[ch-07]{第 7 章}
  \item
    moving average filter：移动平均滤波器的递推实现 --- 282-284
    \dspsym{→} \hyperref[ch-15]{第 15 章}
  \end{itemize}
\item
  \textbf{Reed-Solomon coding} 里德-所罗门编码 --- 361 \dspsym{→}
  \hyperref[ch-22]{第 22 章}
\item
  \textbf{Region-of-convergence} 收敛域 --- 592 \dspsym{→}
  \hyperref[ch-32]{第 32 章}
\item
  \textbf{Resolution} 分辨率

  \begin{itemize}
  \tightlist
  \item
    frequency domain：频域分辨率 --- 173-176 \dspsym{→}
    \hyperref[ch-09]{第 9 章}
  \item
    limiting resolution, images：极限分辨率，图像 --- 426 \dspsym{→}
    \hyperref[ch-25]{第 25 章}
  \item
    spatial domain：空间域分辨率 --- 423-432 \dspsym{→}
    \hyperref[ch-25]{第 25 章}
  \end{itemize}
\item
  \textbf{Retina of the eye} 眼睛的视网膜 --- 378-381 \dspsym{→}
  \hyperref[ch-23]{第 23 章}
\item
  \textbf{Ringing, step response} 振铃，阶跃响应 --- 见
  \textbf{Filters}（滤波器）条下，见 \textbf{Gibbs effect}（吉布斯效应）
\item
  \textbf{RISC} RISC（精简指令集计算机） --- 84 \dspsym{→}
  \hyperref[ch-04]{第 4 章}
\item
  \textbf{RLC circuits} RLC 电路 --- 见 \textbf{Electric circuit
  analysis}（电路分析）
\item
  \textbf{ROC curve} ROC 曲线 --- 见 \textbf{Receiver operating
  characteristic}（接收者操作特性）
\item
  \textbf{Rods and cones in the eye} 眼睛中的视杆细胞与视锥细胞 ---
  379-381 \dspsym{→} \hyperref[ch-23]{第 23 章}
\item
  \textbf{Roll-off} 滚降 --- 见 \textbf{Filters,
  digital}（滤波器，数字）条下
\item
  \textbf{Root-mean-square (rms)} 均方根 --- 14 \dspsym{→}
  \hyperref[ch-02]{第 2 章}
\item
  \textbf{Round-off error} 舍入误差 --- 73-76 \dspsym{→}
  \hyperref[ch-04]{第 4 章}；238 \dspsym{→} \hyperref[ch-12]{第 12
  章}；284 \dspsym{→} \hyperref[ch-15]{第 15 章}；318 \dspsym{→}
  \hyperref[ch-18]{第 18 章}；332 \dspsym{→} \hyperref[ch-19]{第 19 章}
\item
  \textbf{Row major order} 行优先顺序 --- 384 \dspsym{→}
  \hyperref[ch-23]{第 23 章}
\item
  \textbf{Run-length encoding, data compression} 游程编码，数据压缩 ---
  483-484, 500 \dspsym{→} \hyperref[ch-27]{第 27 章}
\item
  \textbf{Running sum} 累加和 --- 126-128 \dspsym{→} \hyperref[ch-07]{第
  7 章}；263 \dspsym{→} \hyperref[ch-14]{第 14 章}
\item
  \textbf{Reverberation in music} 音乐中的混响 --- 5 \dspsym{→}
  \hyperref[ch-01]{第 1 章}
\end{itemize}

\section{S}

\begin{itemize}
\tightlist
\item
  \textbf{s-domain} s 域 --- 581-587 \dspsym{→} \hyperref[ch-32]{第 32
  章}
\item
  \textbf{s-plane} s 平面 --- 581-587 \dspsym{→} \hyperref[ch-32]{第 32
  章}
\item
  \textbf{Saccades of the eye} 眼睛的扫视 --- 381 \dspsym{→}
  \hyperref[ch-23]{第 23 章}
\item
  \textbf{Sampling} 采样 --- 35-44 \dspsym{→} \hyperref[ch-03]{第 3 章}
\item
  \textbf{Sampling aperture in images} 图像的采样孔径 --- 376 \dspsym{→}
  \hyperref[ch-23]{第 23 章}；423, 430-432 \dspsym{→}
  \hyperref[ch-25]{第 25 章}
\item
  \textbf{Sample spacing in images} 图像的采样间距 --- 375-376
  \dspsym{→} \hyperref[ch-23]{第 23 章}；423, 430-432 \dspsym{→}
  \hyperref[ch-25]{第 25 章}
\item
  \textbf{Sampling theorem} 采样定理 --- 39-45 \dspsym{→}
  \hyperref[ch-03]{第 3 章}；430-432 \dspsym{→} \hyperref[ch-25]{第 25
  章}
\item
  \textbf{Scanning probe microscope} 扫描探针显微镜 --- 387 \dspsym{→}
  \hyperref[ch-23]{第 23 章}
\item
  \textbf{SECAM television} SECAM 电视 --- 386 \dspsym{→}
  \hyperref[ch-23]{第 23 章}
\item
  \textbf{Separable image} 可分离图像 --- 404-406 \dspsym{→}
  \hyperref[ch-24]{第 24 章}
\item
  \textbf{Seismology} 地震学 --- 1, 7, 8 \dspsym{→} \hyperref[ch-01]{第
  1 章}；451 \dspsym{→} \hyperref[ch-26]{第 26 章}
\item
  \textbf{Sharpening, image} 锐化，图像 --- 403-404 \dspsym{→}
  \hyperref[ch-24]{第 24 章}
\item
  \textbf{Shift and subtract} 平移相减 --- 402-403 \dspsym{→}
  \hyperref[ch-24]{第 24 章}
\item
  \textbf{Shift invariance} 平移不变性 --- 89-92 \dspsym{→}
  \hyperref[ch-05]{第 5 章}；108 \dspsym{→} \hyperref[ch-06]{第 6 章}
\item
  \textbf{Sidebands in ADC} ADC 中的边带 --- 44-45 \dspsym{→}
  \hyperref[ch-03]{第 3 章}
\item
  \textbf{Sidebands in AM} 调幅中的边带 --- 见 \textbf{Amplitude
  modulation}（幅度调制（调幅））
\item
  \textbf{Sigmoid function} sigmoid 函数 --- 461-463, 471-472 \dspsym{→}
  \hyperref[ch-26]{第 26 章}
\item
  \textbf{Sign and magnitude} 符号-幅值表示 --- 68-69 \dspsym{→}
  \hyperref[ch-04]{第 4 章}
\item
  \textbf{Sign bit} 符号位 --- 68-69 \dspsym{→} \hyperref[ch-04]{第 4
  章}
\item
  \textbf{Signed fraction} 有符号小数 --- 514 \dspsym{→}
  \hyperref[ch-28]{第 28 章}
\item
  \textbf{Signal (defined)} 信号（定义） --- 11 \dspsym{→}
  \hyperref[ch-02]{第 2 章}；87 \dspsym{→} \hyperref[ch-05]{第 5 章}
\item
  \textbf{Signal-to-noise ratio} 信噪比 --- 17 \dspsym{→}
  \hyperref[ch-02]{第 2 章}；432-436 \dspsym{→} \hyperref[ch-25]{第 25
  章}
\item
  \textbf{Sinc function} sinc 函数

  \begin{itemize}
  \tightlist
  \item
    aliasing of：其混叠 --- 212-214 \dspsym{→} \hyperref[ch-11]{第 11
    章}
  \item
    DAC reconstruction：DAC 重建 --- 46-47 \dspsym{→}
    \hyperref[ch-03]{第 3 章}
  \item
    Fourier transform of rectangular pulse：矩形脉冲的傅里叶变换 ---
    212-215 \dspsym{→} \hyperref[ch-11]{第 11 章}；285-287 \dspsym{→}
    \hyperref[ch-16]{第 16 章}
  \item
    Two-dimensional (image)：二维（图像） --- 400-401 \dspsym{→}
    \hyperref[ch-24]{第 24 章}
  \item
    Windowed-sinc filter：加窗 sinc 滤波器 --- 283-285 \dspsym{→}
    \hyperref[ch-15]{第 15 章}、\hyperref[ch-16]{第 16 章}
  \end{itemize}
\item
  \textbf{Single precision} 单精度 --- 70-76 \dspsym{→}
  \hyperref[ch-04]{第 4 章}
\item
  \textbf{Sinusoidal fidelity} 正弦保真性 --- 见
  \textbf{Linearity}（线性）
\item
  \textbf{SIRT} SIRT（联立迭代重建技术） --- 442 \dspsym{→}
  \hyperref[ch-25]{第 25 章}
\item
  \textbf{Skeletonization of binary images} 二值图像的骨架化 --- 438-442
  \dspsym{→} \hyperref[ch-25]{第 25 章}
\item
  \textbf{Smoothing filter} 平滑滤波器 --- 见
  \textbf{Filters}（滤波器）条下
\item
  \textbf{Space exploration} 太空探索 --- 见
  \textbf{Astrophotography}（天文摄影）
\item
  \textbf{Spectrogram of speech} 语音的频谱图 --- 365-367 \dspsym{→}
  \hyperref[ch-22]{第 22 章}
\item
  \textbf{Speed} 速度

  \begin{itemize}
  \tightlist
  \item
    convolution vs.~FFT convolution：卷积与 FFT 卷积对比 --- 318
    \dspsym{→} \hyperref[ch-18]{第 18 章}
  \item
    FIR vs.~IIR filters：FIR 与 IIR 滤波器对比 --- 346-350 \dspsym{→}
    \hyperref[ch-21]{第 21 章}
  \item
    hardware：硬件 --- 80-84 \dspsym{→} \hyperref[ch-04]{第 4 章}
  \item
    image convolution：图像卷积 --- 421-422 \dspsym{→}
    \hyperref[ch-24]{第 24 章}
  \item
    program language：编程语言 --- 76-80 \dspsym{→} \hyperref[ch-04]{第
    4 章}
  \item
    programming style：编程风格 --- 84-86 \dspsym{→} \hyperref[ch-04]{第
    4 章}
  \end{itemize}
\item
  \textbf{Spatial domain (defined)} 空间域（定义） --- 12 \dspsym{→}
  \hyperref[ch-02]{第 2 章}；373 \dspsym{→} \hyperref[ch-23]{第 23 章}
\item
  \textbf{Spatial resolution} 空间分辨率 --- 见
  \textbf{Resolution}（分辨率）条下
\item
  \textbf{Speak \& spell} Speak \& Spell（德州仪器的拼读学习机） --- 6
  \dspsym{→} \hyperref[ch-01]{第 1 章}；366 \dspsym{→}
  \hyperref[ch-22]{第 22 章}
\item
  \textbf{Spectral analysis} 频谱分析 --- 169-176 \dspsym{→}
  \hyperref[ch-09]{第 9 章}
\item
  \textbf{Spectral inversion} 频谱反转 --- 见
  \textbf{Filters}（滤波器）条下
\item
  \textbf{Spectral response of the eye} 眼睛的光谱响应 --- 380-381
  \dspsym{→} \hyperref[ch-23]{第 23 章}
\item
  \textbf{Spectral reversal} 频谱翻转 --- 见 \textbf{Filters,
  digital}（滤波器，数字）条下
\item
  \textbf{Speech} 语音

  \begin{itemize}
  \tightlist
  \item
    digital recording：数字录音 --- 6 \dspsym{→} \hyperref[ch-01]{第 1
    章}
  \item
    generation：语音生成 --- 6 \dspsym{→} \hyperref[ch-01]{第 1
    章}；357, 364-368 \dspsym{→} \hyperref[ch-22]{第 22 章}
  \item
    recognition：语音识别 --- 6-7 \dspsym{→} \hyperref[ch-01]{第 1
    章}；364-368 \dspsym{→} \hyperref[ch-22]{第 22 章}
  \item
    speak \& spell：Speak \& Spell 拼读学习机 --- 6 \dspsym{→}
    \hyperref[ch-01]{第 1 章}
  \item
    vocal tract simulation：声道模拟 --- 6 \dspsym{→}
    \hyperref[ch-01]{第 1 章}
  \end{itemize}
\item
  \textbf{Sonar} 声呐 --- 1, 7-8 \dspsym{→} \hyperref[ch-01]{第 1
  章}；32-34 \dspsym{→} \hyperref[ch-02]{第 2 章}；88 \dspsym{→}
  \hyperref[ch-05]{第 5 章}；171-173 \dspsym{→} \hyperref[ch-09]{第 9
  章}；222-224 \dspsym{→} \hyperref[ch-11]{第 11 章}
\item
  \textbf{Square PSF} 方形点扩散函数 --- 400-401 \dspsym{→}
  \hyperref[ch-24]{第 24 章}
\item
  \textbf{Stability} 稳定性 --- 339 \dspsym{→} \hyperref[ch-20]{第 20
  章}；541 \dspsym{→} \hyperref[ch-29]{第 29
  章}（疑为旧版页码，按内容应在 589 \dspsym{→} \hyperref[ch-32]{第 32
  章}）；600-601 \dspsym{→} \hyperref[ch-32]{第 32 章}；609 \dspsym{→}
  \hyperref[ch-33]{第 33 章}
\item
  \textbf{Standard deviation} 标准差 --- 13-17, 20-22 \dspsym{→}
  \hyperref[ch-02]{第 2 章}；434-435 \dspsym{→} \hyperref[ch-25]{第 25
  章}
\item
  \textbf{Stationary} 平稳 --- 19 \dspsym{→} \hyperref[ch-02]{第 2 章}
\item
  \textbf{Static linearity} 静态线性 --- 92-93 \dspsym{→}
  \hyperref[ch-05]{第 5 章}
\item
  \textbf{Statistical variation} 统计变异 --- 见 \textbf{Noise}（噪声）
\item
  \textbf{Steepest decent, neural network learning}
  最速下降，神经网络学习 --- 470-473 \dspsym{→} \hyperref[ch-26]{第 26
  章}
\item
  \textbf{Step decomposition} 阶跃分解 --- 见
  \textbf{Decomposition}（分解）条下
\item
  \textbf{Step response} 阶跃响应 --- 见 \textbf{Filters}（滤波器）条下
\item
  \textbf{Stereo audio} 立体声音频 --- 362 \dspsym{→}
  \hyperref[ch-22]{第 22 章}
\item
  \textbf{Stopband (defined)} 阻带（定义） --- 268 \dspsym{→}
  \hyperref[ch-14]{第 14 章}
\item
  \textbf{Stretch, grayscale} 拉伸，灰度 --- 390-391 \dspsym{→}
  \hyperref[ch-23]{第 23 章}
\item
  \textbf{Strong law of large numbers} 强大数定律 --- 18 \dspsym{→}
  \hyperref[ch-02]{第 2 章}
\item
  \textbf{Subpixel interpolation} 亚像素插值 --- 395-396 \dspsym{→}
  \hyperref[ch-23]{第 23 章}
\item
  \textbf{Substitution, using complex numbers by} 代换，用复数做代换 ---
  557-559 \dspsym{→} \hyperref[ch-30]{第 30 章}
\item
  \textbf{Superposition} 叠加 --- 见 \textbf{Linearity}（线性）条下
\item
  \textbf{Symmetry, left-right} 对称，左右 --- 见 \textbf{Phase} 条下的
  zero phase
\item
  \textbf{Synthesis} 合成 --- 见
  \textbf{Linearity}（线性）；\textbf{Fourier transform, Discrete
  Fourier transform}条下
\item
  \textbf{System (defined)} 系统（定义） --- 87 \dspsym{→}
  \hyperref[ch-05]{第 5 章}
\item
  \textbf{Systematic errors} 系统误差 --- 32-34 \dspsym{→}
  \hyperref[ch-02]{第 2 章}
\end{itemize}

\section{T}

\begin{itemize}
\tightlist
\item
  \textbf{t-carrier system} T 载波系统 --- 4 \dspsym{→}
  \hyperref[ch-01]{第 1 章}
\item
  \textbf{Target detection} 目标检测 --- 451-458 \dspsym{→}
  \hyperref[ch-26]{第 26 章}
\item
  \textbf{Tchebysheff} Tchebysheff（切比雪夫的另一种拼法） --- 见
  \textbf{Chebyshev}
\item
  \textbf{Telecommunications} 电信 --- 4-5 \dspsym{→}
  \hyperref[ch-01]{第 1 章}
\item
  \textbf{Television} 电视 --- 206 \dspsym{→} \hyperref[ch-10]{第 10
  章}；374, 384-386 \dspsym{→} \hyperref[ch-23]{第 23 章}；501
  \dspsym{→} \hyperref[ch-27]{第 27 章}
\item
  \textbf{Text recognition, neural network} 文字识别，神经网络 ---
  465-476 \dspsym{→} \hyperref[ch-26]{第 26 章}
\item
  \textbf{TIFF image, data compression} TIFF 图像，数据压缩 --- 488
  \dspsym{→} \hyperref[ch-27]{第 27 章}
\item
  \textbf{Timbre} 音色 --- 355-358 \dspsym{→} \hyperref[ch-22]{第 22 章}
\item
  \textbf{Time domain (defined)} 时域（定义） --- 12 \dspsym{→}
  \hyperref[ch-02]{第 2 章}；146-147 \dspsym{→} \hyperref[ch-08]{第 8
  章}
\item
  \textbf{Time domain encoding} 时域编码 --- 见
  \textbf{Information}（信息）条下
\item
  \textbf{Threshold} 阈值 --- 453-458 \dspsym{→} \hyperref[ch-26]{第 26
  章}
\item
  \textbf{Transfer function} 传递函数 --- 594 \dspsym{→}
  \hyperref[ch-32]{第 32 章}；611-612 \dspsym{→} \hyperref[ch-33]{第 33
  章}
\item
  \textbf{Transform (defined)} 变换（定义） --- 146 \dspsym{→}
  \hyperref[ch-08]{第 8 章}
\item
  \textbf{Transform data compression} 变换数据压缩 --- 494-500
  \dspsym{→} \hyperref[ch-27]{第 27 章}
\item
  \textbf{Triangular pulse} 三角形脉冲 --- 216-217 \dspsym{→}
  \hyperref[ch-11]{第 11 章}；281 \dspsym{→} \hyperref[ch-15]{第 15 章}
\item
  \textbf{Trigonometric functions} 三角函数 --- 85-86 \dspsym{→}
  \hyperref[ch-04]{第 4 章}；165-166 \dspsym{→} \hyperref[ch-08]{第 8
  章}
\item
  \textbf{True-positive (true-negative)} 真阳性（真阴性） --- 453-454
  \dspsym{→} \hyperref[ch-26]{第 26 章}
\item
  \textbf{Tschebyscheff} Tschebyscheff（切比雪夫的另一种拼法） --- 见
  \textbf{Chebyshev}
\item
  \textbf{Two's complement} 补码 --- 69-70 \dspsym{→}
  \hyperref[ch-04]{第 4 章}
\end{itemize}

\section{U}

\begin{itemize}
\tightlist
\item
  \textbf{µ255 law companding} μ255 律压扩 --- 362-364 \dspsym{→}
  \hyperref[ch-22]{第 22 章}
\item
  \textbf{Unit circle, z-plane} 单位圆，z 平面 --- 608-609 \dspsym{→}
  \hyperref[ch-33]{第 33 章}
\item
  \textbf{Unsigned integer} 无符号整数 --- 68-69 \dspsym{→}
  \hyperref[ch-04]{第 4 章}
\end{itemize}

\section{V}

\begin{itemize}
\tightlist
\item
  \textbf{Variance} 方差 --- 14 \dspsym{→} \hyperref[ch-02]{第 2 章}
\item
  \textbf{Von Neumann architecture} 冯·诺依曼结构 --- 509 \dspsym{→}
  \hyperref[ch-28]{第 28 章}
\item
  \textbf{Voiced sound in speech} 语音中的浊音 --- 364 \dspsym{→}
  \hyperref[ch-22]{第 22 章}
\item
  \textbf{Voiceprint of speech} 语音的声纹 --- 365-367 \dspsym{→}
  \hyperref[ch-22]{第 22 章}
\end{itemize}

\section{W}

\begin{itemize}
\tightlist
\item
  \textbf{Warping images} 图像变形 --- 394-396 \dspsym{→}
  \hyperref[ch-23]{第 23 章}
\item
  \textbf{Well, in CCDs} 阱，CCD 中 --- 382-384 \dspsym{→}
  \hyperref[ch-23]{第 23 章}
\item
  \textbf{Wiener filter} 维纳滤波器 --- 见 \textbf{Filters,
  digital}（滤波器，数字）条下
\item
  \textbf{Windows} 窗函数

  \begin{itemize}
  \tightlist
  \item
    Bartlett window：巴特利特窗 --- 288 \dspsym{→} \hyperref[ch-16]{第
    16 章}
  \item
    Blackman：布莱克曼窗 --- 175-176 \dspsym{→} \hyperref[ch-09]{第 9
    章}；282-283 \dspsym{→} \hyperref[ch-15]{第 15 章}；288 \dspsym{→}
    \hyperref[ch-16]{第 16 章}
  \item
    Hamming：汉明窗 --- 170-171, 174-176 \dspsym{→} \hyperref[ch-09]{第
    9 章}；282-283 \dspsym{→} \hyperref[ch-15]{第 15 章}；288 \dspsym{→}
    \hyperref[ch-16]{第 16 章}
  \item
    Hanning：汉宁窗 --- 288 \dspsym{→} \hyperref[ch-16]{第 16 章}
  \item
    raised cosine：升余弦窗 --- 288 \dspsym{→} \hyperref[ch-16]{第 16
    章}
  \item
    rectangular：矩形窗 --- 175-176 \dspsym{→} \hyperref[ch-09]{第 9
    章}；212-215 \dspsym{→} \hyperref[ch-11]{第 11 章}；288 \dspsym{→}
    \hyperref[ch-16]{第 16 章}
  \item
    in spectral analysis：在频谱分析中 --- 169-176, 173-176 \dspsym{→}
    \hyperref[ch-09]{第 9 章}
  \end{itemize}
\end{itemize}

\section{X}

\begin{itemize}
\tightlist
\item
  \textbf{X-ray imaging} X 射线成像

  \begin{itemize}
  \tightlist
  \item
    另见 \textbf{computed tomography}（计算机断层成像）
  \item
    airport baggage scanner：机场行李扫描仪 --- 402-403 \dspsym{→}
    \hyperref[ch-24]{第 24 章}
  \item
    detection by phosphor layer：用荧光层探测 --- 423-424 \dspsym{→}
    \hyperref[ch-25]{第 25 章}
  \item
    DSP improvements to：DSP 带来的改进 --- 9 \dspsym{→}
    \hyperref[ch-01]{第 1 章}
  \item
    measuring MTF：测量 MTF --- 425 \dspsym{→} \hyperref[ch-25]{第 25
    章}
  \item
    image noise of：其图像噪声 --- 434-436 \dspsym{→}
    \hyperref[ch-25]{第 25 章}
  \end{itemize}
\end{itemize}

\section{Z}

\begin{itemize}
\tightlist
\item
  \textbf{z-domain} z 域 --- 605-610 \dspsym{→} \hyperref[ch-33]{第 33
  章}
\item
  \textbf{z-plane} z 平面 --- 605-610 \dspsym{→} \hyperref[ch-33]{第 33
  章}
\item
  \textbf{z-transform} z 变换 --- 334-335 \dspsym{→} \hyperref[ch-20]{第
  20 章}；605-610 \dspsym{→} \hyperref[ch-33]{第 33 章}
\item
  \textbf{Zero phase} 零相位 --- 见 \textbf{Phase}（相位）条下
\item
  \textbf{Zeros} 零点 --- 见 \textbf{Poles and zeros}（极点与零点）
\item
  \textbf{Zeroth-order-hold, DAC} 零阶保持，DAC --- 46-47 \dspsym{→}
  \hyperref[ch-03]{第 3 章}
\end{itemize}

\begin{center}\rule{0.5\linewidth}{0.5pt}\end{center}

\section{译注：原书索引中的问题}

\begin{enumerate}
\def\labelenumi{\arabic{enumi}.}
\tightlist
\item
  \textbf{疑为旧版页码的几处。} 下面这些页码落在第 28、29 章（讲 DSP
  芯片），但那几页里并没有对应内容；按内容查，它们应在第 30～32
  章，大多正好往后差 48 页，也就是第 28、29
  两章的篇幅。推测这几条沿用了加入这两章之前的旧版页码。条目里已在原页码后注明按内容核对出的位置：

  \begin{itemize}
  \tightlist
  \item
    Discrete Fourier transform 下的 complex DFT 519-532、real DFT
    519-521 \dspsym{→} 应为 567-580、567-569（第 31 章）；
  \item
    Fourier transform 下的 scaling 529-531 \dspsym{→} 应为 577-579（第
    31 章）；
  \item
    Phasor transform 515 \dspsym{→} 应为 563（第 30 章）；
  \item
    Filters, analog 下的 notch filter 519-520 \dspsym{→}
    陷波滤波器实际在 565-566（第 30 章）；
  \item
    Stability 与 Filters, analog 下的 stability 541 \dspsym{→} 应为
    589（第 32 章）。
  \end{itemize}
\item
  Complex numbers 下的 commutative property
  印成「5054」，对照渲染页确认原书就是这样印的，应为 554（第 30 章）。
\item
  Convolution 下的 execution time
  印成「318-422」，跨了七章，显然有误；按内容推测应为 318（第 18 章）与
  421-422（第 24 章），条目里照原书保留，只按两端映射。
\item
  Filters, recursive 下的 stability 写作「339, 599, 598-599, 609」，599
  与 598-599 重复，照原样保留。
\item
  Fourier transform 下的「discrete time Fourier series (DTFT)」应为
  discrete time Fourier transform（离散时间傅里叶变换）。
\item
  拼写错误，照原样保留英文：Orthognal（应为 Orthogonal）、Format
  frequencies（应为 Formant，共振峰）、Steepest decent（应为
  descent）。Dilatation、Tchebysheff、Tschebyscheff
  是异体拼写，原书本来就是作为参见条目列出的。
\item
  原书索引奇数页的页眉印成了「Glossary」（应为 Index）。
\item
  有几处没有严格按字母排序，例如 Caruso 排在 Cascaded 后、Hyperspace
  排在 Huffman 前、Nyquist 排在 Nuclear 前、Sonar 排在 Speech
  后、Threshold 排在 Time domain encoding 后、Reverberation in music
  排在 R 部末尾。这里照原书顺序保留。
\end{enumerate}
